\section{When chart-wise gluing fails, and what survives}
\label{sec:gluing:obstructions}

The chart-wise assemblies of \S\S\ref{sec:gluing:toric-cech}--\ref{sec:gluing:k3e} glue cohomological Hall algebras from local data by cocycles valued in explicitly-named categories. Four features of the geometry power that machinery: a toric atlas cover by $\bC^3$-pieces, vanishing of the BCOV one-loop anomaly on each piece, torus-equivariant localisation on the moduli of simples, and a finite quiver with potential encoding the D-brane category. The present section demonstrates that each feature fails on a generic compact Calabi--Yau threefold, derives the four failures from first principles, and tracks what is and is not recovered by stratum-independent machinery.

The four failures are not independent. By Sumihiro 1974 Thm.~1 the pair $(\mathbf{O_1}) \Leftrightarrow (\mathbf{O_3})$ is a single toric-rigidity biconditional; by Van den Bergh + Bocklandt + Hartogs that biconditional implies $(\mathbf{O_4})$; and $(\mathbf{O_2})$ is the unique Hodge-theoretic obstruction (Costello--Li \texttt{arXiv:1606.00365} Prop.~5.2), logically independent of the toric-rigidity axis. Thus the four failures reduce to a toric-rigidity axis, a BCOV axis, and one Hodge-theoretic independence statement. The remainder of the section records what Davison integrality, Kontsevich--Soibelman motivic wall-crossing, and Borcherds automorphic forms recover on compact CY$_3$, and what they do not, concluding with the five-fold catalogue of obstructions to a global NCCR on $K3\times E$ and the Serre-equivariant quasi-NCCR substitute.
\subsection{The four obstructions}
\label{sec:four-obstructions}

Fix a smooth projective Calabi--Yau threefold $X$ over $\bC$. We wish
to assemble a cohomological Hall algebra $\CoHA(X)$ from chart-wise
local pieces, as in Parts~II--V. Four features of the geometry are
required.

\begin{itemize}[leftmargin=2em]
\item[$(\mathbf{C_1})$] A \emph{toric atlas}: an open cover $\{U_\alpha\}$
 of $X$ with each $U_\alpha$ biholomorphic to $\bC^3$, intersections
 $U_\alpha \cap U_\beta$ toric, and transition cocycles in
 $\mathrm{GL}_3(\bC)$.
\item[$(\mathbf{C_2})$] \emph{BCOV-anomaly vanishing}: the Costello--Li
 one-loop cocycle $\alpha_{\mathrm{BCOV}} = \tfrac{\chi(X)}{24}
 \cdot \mathrm{tr}\,\mathrm{At}(T_X) \in H^{1}(X, \cO_X) =
 H^{0,1}(X, \cO_X)$ vanishes, allowing a factorisation-algebra
 quantisation of 5D hCS on $X$.
\item[$(\mathbf{C_3})$] \emph{Torus equivariance}: an algebraic torus
 $T \curvearrowright X$ with finitely many fixed points on each
 moduli space $\cM^n(X)$ of BPS $n$-objects, so that
 Atiyah--Bott localisation reduces enumerative invariants to fixed-point
 sums.
\item[$(\mathbf{C_4})$] \emph{Finite quiver with potential}: a tilting
 object $T \in D^b(\Coh\,X)$ whose endomorphism algebra $A = \End(T)$
 is a Jacobi algebra of a finite quiver $Q$ with superpotential $W$.
\end{itemize}

On the resolved conifold all four hold; on local $\bP^2$ all four hold;
on any toric CY$_3$ without compact $4$-cycles all four hold. The
first-principles question of this part is: which features survive on a
compact CY$_3$?

\begin{theorem}[Four obstructions]
\label{thm:four-obstructions}
Let $X$ be a smooth projective Calabi--Yau threefold with $h^{2,0}(X) = 0$
and $\Pic^0(X) = 0$. Each of the four features $(\mathbf{C_1})$ through
$(\mathbf{C_4})$ fails:
\begin{enumerate}[label=\emph{$(\mathbf{O_\arabic*})$}, leftmargin=2.5em]
\item $X$ admits no toric atlas;
\item the BCOV one-loop anomaly $\alpha_{\mathrm{BCOV}}$ need not
 vanish;
\item $X$ admits no algebraic torus action with isolated fixed points on
 simples;
\item $D^b(\Coh\,X)$ admits no tilting object with $\End(T)$ a finite
 Jacobi algebra.
\end{enumerate}
The equivalence $(\mathbf{O_1}) \Leftrightarrow (\mathbf{O_3})$ (Sumihiro
1974, Proposition~\ref{prop:sumihiro-biconditional}) and
$(\mathbf{O_1}) \Rightarrow (\mathbf{O_4})$ (Van den Bergh + Bocklandt
+ Hartogs); $(\mathbf{O_2})$ is logically independent (Costello--Li
\texttt{arXiv:1606.00365} Prop.~5.2). The four named features carry
$2.5$ units of independent logical content
(Proposition~\ref{prop:structural-content}).
\end{theorem}

\begin{proof}[Proof of $(\mathbf{O_1})$]
A toric atlas $\{U_\alpha \simeq \bC^3\}$ with overlaps in
$\mathrm{GL}_3(\bC)$-monomial coordinates forces $X$ to be a smooth
projective toric threefold: the implicit torus $(\bC^*)^3 \subset
\bC^3$ in each chart glues via the monomial cocycles to a global dense
open torus $T \subset X$. By the fundamental theorem of smooth toric
varieties (Demazure; Fulton, \emph{Introduction to Toric Varieties}
\S3.4, Thm.~3.1.5), such $X$ corresponds bijectively to a complete
regular fan $\Sigma$ in the lattice $N = \bZ^3$, with the $3$-dimensional
cones of $\Sigma$ indexing the charts and the rays $\rho \in \Sigma(1)$
indexing the toric-divisor generators.

The Calabi--Yau condition $\omega_X \simeq \cO_X$ on a toric variety
translates, via the exact sequence
\[
 0 \longrightarrow M \xrightarrow{\;\psi\;}
 \bigoplus_{\rho \in \Sigma(1)} \bZ\cdot D_\rho
 \longrightarrow \Pic(X) \longrightarrow 0,
\]
into the existence of $\psi \in M = \Hom(N, \bZ)$ with
$\langle \psi, v_\rho\rangle = 1$ for every ray generator
$v_\rho \in N$ (Fulton \S3.4, the CY datum). All ray generators lie on
the affine hyperplane
\[
 H_\psi \;=\; \{v \in N \otimes \bR : \langle \psi, v\rangle = 1\}
 \;\subset\; \bR^3,
\]
a two-dimensional affine plane in $\bR^3$ not passing through the origin.

Completeness (Fulton \S3.4, Thm.~3.1.5) is the condition
$|\Sigma| := \bigcup_{\sigma \in \Sigma} \sigma = N_\bR = \bR^3$: every
$v \in \bR^3$ lies in some cone of $\Sigma$, equivalently is a
non-negative $\bR_{\geq 0}$-linear combination of ray generators
$v_\rho$. If all ray generators lie on the affine hyperplane
$H_\psi = \{v : \langle \psi, v \rangle = 1\}$, then every non-negative
combination $\sum_\rho \lambda_\rho v_\rho$ with $\lambda_\rho \geq 0$
satisfies $\langle \psi, \sum_\rho \lambda_\rho v_\rho \rangle =
\sum_\rho \lambda_\rho \geq 0$, placing $|\Sigma|$ in the closed
half-space $\{v : \langle \psi, v\rangle \geq 0\}$, strictly contained
in $\bR^3$. Hence no vector $v$ with $\langle \psi, v\rangle < 0$ is
covered by the fan: $|\Sigma| \neq \bR^3$, contradicting completeness.
No complete smooth toric CY$_3$ exists (Fulton \S3.4, Thm.~3.1.5 and
the CY datum \S3.4; equivalently, projective toric CY$_d$ for
$d \geq 1$ are the Gorenstein toric varieties associated to reflexive
polytopes; singular, \emph{never} smooth).

The argument is entirely local-to-global: any compact CY$_3$ admitting
a $\bC^3$-chart atlas with monomial transition cocycles would be a
smooth projective toric CY$_3$, and no such object exists. The quintic
$X_5 \subset \bP^4$, the bicubic $X_{3,3}$, the $(2,4)$-complete
intersection, the Schoen threefold, $K3 \times E$, and every other
compact CY$_3$ in the physicist census fail $(\mathbf{C_1})$ by the
same structural reason.
\end{proof}

\begin{proof}[Proof of $(\mathbf{O_2})$]
Costello--Li (\texttt{arXiv:1606.00365} Prop.~5.2; cf.\
\texttt{arXiv:1605.09656}) identify the one-loop obstruction to
quantising the holomorphic BCOV factorisation algebra on a
Calabi--Yau threefold $X$ with the \emph{Atiyah--dilaton cocycle}
\[
 \alpha_{\mathrm{BCOV}}(X)
 \;=\; \frac{\chi_{\mathrm{top}}(X)}{24}\;\tr\,\mathrm{At}(T_X)
 \;\in\; H^{1}(X,\cO_X),
\]
a class in the \emph{structure-sheaf} first cohomology, not in any
symmetric power of $T_X^\ast$. The Atiyah class
$\mathrm{At}(T_X)\in H^1(X,\Omega^1_X\otimes\End T_X)$ of the
holomorphic tangent bundle has fibrewise trace
$\tr\,\mathrm{At}(T_X)\in H^1(X,\Omega^1_X)$ equal to
$\mathrm{At}(\det T_X)=-\mathrm{At}(\omega_X)$; on a Calabi--Yau the
chosen holomorphic volume form kills $\mathrm{At}(\omega_X)$ in
$H^1(X,\Omega^1_X)$, and the CY residue contraction
$\Omega^1_X\otimes\omega_X\simeq\Omega^1_X\to\cO_X$ transports
$\tr\,\mathrm{At}(T_X)$ onto its dilaton zero-mode in
$H^1(X,\cO_X)=H^{0,1}(X,\cO_X)$; a single cohomology group, one
scalar obstruction per $X$, read off the Atiyah--dilaton pairing.
Primary sources: Costello--Li \texttt{arXiv:1606.00365} Prop.~5.2
for the identification; BCOV \texttt{arXiv:hep-th/9309140}
Eq.~(5.11) for the $\chi_{\mathrm{top}}/24$ prefactor via
Chern--Gauss--Bonnet on the one-loop trace diagram; Polyakov
\emph{Gauge Fields and Strings} Ch.~9 for the two-dimensional
conformal-anomaly ancestor from which the chiral-BCOV descends;
Kapranov \texttt{arXiv:math/9811023} for the identification of
$\tr\,\mathrm{At}$ with the dilaton one-form on a Calabi--Yau.

The prefactor $\chi(X)/24 \in \Q$ is the universal
Chern--Gauss--Bonnet coefficient of the one-loop holomorphic-BCOV
trace diagram. Its derivation: the 5D hCS one-loop partition function
on $\R \times X$ reduces via Duhamel to the integrated
$\widehat A$-genus twist of the identity trace on $T_X$, giving
$\int_X \mathrm{td}(T_X) \mathrm{ch}(T_X) = \chi(X)$ by
Chern--Gauss--Bonnet on the target threefold, with the universal
normalisation $1/24$ dictated by the BV measure on the global ghost
tower (Costello \texttt{arXiv:1605.09656} \S 12.3).

Throughout this section we write $[\Omega_X]^{0,1}$ as shorthand for
the Dolbeault representative of $\mathrm{tr}\,\mathrm{At}(T_X)$ inside
$H^{0,1}(X, \cO_X)$, obtained by projecting $\mathrm{tr}\,\mathrm{At}(T_X)
\in H^1(X, \Omega^1_X)$ along the CY residue contraction
$\Omega^1_X \otimes \omega_X \simeq \Omega^1_X \to \cO_X$; the two
notations refer to the same class, distinguished only by whether one
reads the trace through the volume form $\Omega_X \in
H^0(X, \omega_X)$ or through the tangent-bundle Atiyah class.

\emph{Vanishing on non-compact toric CY$_3$.} On the resolved conifold
$\bY = \mathrm{Tot}(\cO(-1)\oplus\cO(-1) \to\bP^1)$ two orthogonal
mechanisms collapse $\alpha_{\mathrm{BCOV}}(\bY)$ to zero. First, the
retraction $\bY \xrightarrow{r} \bP^1$ induces the isomorphism
$r^\ast: H^{0,1}(\bP^1, \cO_{\bP^1}) \xrightarrow{\sim}
H^{0,1}(\bY, \cO_\bY)$, with target
$H^{0,1}(\bP^1, \cO_{\bP^1}) = H^1(\bP^1, \cO_{\bP^1}) = 0$
(Serre duality: $H^0(\bP^1, \cO(-2)) = 0$). Second, independently,
$\Omega_\bY$ descends from the deformed conifold $\{xy - zw = 0\}$
via Atiyah flop as a trivialisation of $\omega_\bY$ with no
$(0,1)$-component at infinity; the Atiyah class of a trivial line
bundle vanishes. The two mechanisms agree:
$\alpha_{\mathrm{BCOV}}(\bY) = 0$ as a class in $H^{0,1}(\bY, \cO_\bY)$;
cross-ref Remark~\ref{rem:anomaly-vanishing} and
Lemma~\ref{lem:chi-top-bcov-split}. The topological Euler
characteristic $\chi_{\mathrm{top}}(\bY) = \chi(\bP^1) = 2$ is finite
and non-zero; it is the \emph{cohomology} factor $[\Omega_\bY]^{0,1}$
that vanishes, not the topological scalar. On a toric CY$_3$ without
compact $4$-cycles (local $\bP^2$ excluded, since it has a compact
divisor class), the same retraction argument onto the $1$-skeleton of
the toric fan applies.

\emph{The $a_{\mathrm{BV}}$ curving on the conifold.}
Although $\alpha_{\mathrm{BCOV}}(\bY)=0$ as a bulk class in
$H^1(\bY,\cO_{\bY})$, the boundary-link $T^{1,1}\simeq S^2\times S^3$
carries a residual ghost-mode curving. The one-loop BCOV diagram on
the cone, when regulated by the heat kernel on $L^5=T^{1,1}$,
produces a $\zeta(0)$-regularised residue of the $\beta\gamma$-ghost
central charge $c=-2$ (Kausch \texttt{arXiv:hep-th/9510149}
Eq.~(4.12)), and the boundary integral over $T^{1,1}$ fixes
\[
 a_{\mathrm{BV}}(\bY) \;=\; -1,
\]
the exact value dictated by the ghost-mode balance with the
fermionic orientation on the $\cO(-1)\oplus\cO(-1)$ bundle. This
is the Polyakov boundary computation of the conifold BCOV
residue; the Costello--Li propagator on the $\bP^1$-retract agrees.
The coefficient $a_{\mathrm{BV}}$ is \emph{distinct} from
$\kappa_{\mathrm{ch}}$: the former is the $\zeta$-regularised
ghost-mode residue of the boundary link, the latter is the
chiral-supertrace shadow of the bulk CY category. On the conifold,
$\kappa_{\mathrm{ch}}^{\mathrm{DT}}(\bY)=+1$ from the Szendr\H{o}i
two-chart DT generating function, and the ghost-mode balance
$\kappa_{\mathrm{ch}}^{\mathrm{DT}}(\bY)+
a_{\mathrm{BV}}(\bY)=0=\kappa_{\mathrm{cat}}(\bY)$ holds by
direct Hodge-supertrace on the Kähler deformation-retract $\bP^1$,
which is the class-G free-field manifestation of the balance law
stated in Remark~\ref{rem:kappa-ch-bv-distinction}. This balance
is specific to class-G free-field CY$_3$ at $\chi_{\mathrm{top}}=2$;
it is \emph{not} the Theorem~C derived-centre ceiling and does
\emph{not} generalise to other toric CY$_3$.

\emph{Four-case stratification by $\chi_{\mathrm{top}}$ and
$h^{0,1}$.} The Atiyah--dilaton cocycle factors as topological
scalar $\times$ Dolbeault class, and each factor vanishes on
distinct geometric data:
\begin{itemize}[leftmargin=2em]
 \item[$(\beta_1)$] \emph{Non-compact, retractable} (resolved
 conifold $\bY$; any toric CY$_3$ without compact $4$-cycles): bulk
 $H^1(X_{\mathrm{con}},\cO)=0$ because the affine projection
 $\pi$ collapses bulk Dolbeault cohomology onto the $1$-skeleton of
 the toric fan. The leading Atiyah--dilaton class vanishes in the
 bulk by Remark~\ref{rem:anomaly-vanishing}; the topological scalar
 $\chi_{\mathrm{top}}(\bY)=2$ survives as a \emph{boundary} curving
 computed by Sasaki--Einstein-link $\zeta$-regularisation below.
 \item[$(\beta_2)$] \emph{Compact strict CY$_3$} (quintic $X_5$,
 bicubic $X_{3,3}$, Schoen): the strict-CY Hodge condition
 $h^{1,0}=0$ combined with $\omega_X\simeq\cO_X$ forces
 $h^{0,1}(X)=h^{1,0}(X)=0$, so the leading cocycle target $H^1(X,
 \cO_X)=0$ is itself trivial. $\alpha_{\mathrm{BCOV}}(X_5)=0$
 despite $\chi_{\mathrm{top}}(X_5)=-200$; the target vanishes, so
 the class vanishes trivially.
 \item[$(\beta_3)$] \emph{$K3\times E$}: Künneth kills the
 topological scalar, $\chi_{\mathrm{top}}(K3\times E)=\chi(K3)\cdot
 \chi(E)=24\cdot 0=0$; independently, $h^{0,1}(K3\times E)=
 h^{0,1}(E)=1$, so $\tr\,\mathrm{At}(T_X)\neq 0$ in
 $H^1(K3\times E,\cO)$. The product
 $\alpha_{\mathrm{BCOV}}(K3\times E)=0$ vanishes for
 \emph{topological} reasons (Künneth), not cohomological.
 \item[$(\beta_4)$] \emph{Compact CY$_3$ with abelian factor}
 (abelian threefold $A^3$; any CY$_3$ with positive-dimensional
 abelian summand in Bogomolov--Beauville): $\chi_{\mathrm{top}}=0$
 and $h^{0,1}\neq 0$, so the cocycle vanishes topologically while
 the cohomology target is non-trivial; the deformation
 \emph{gradient} of $\alpha_{\mathrm{BCOV}}$ detects the abelian
 direction.
\end{itemize}
None of the four cases places $\alpha_{\mathrm{BCOV}}$ in a
non-zero class of $H^1(X,\cO_X)$ on a compact smooth CY$_3$; the
leading Costello--Li obstruction vanishes on every named witness.
Higher-weight BCOV obstructions on strict CY$_3$ (quintic) occupy
the iterated one-loop BV tower above the leading Atiyah trace, not
a structure-sheaf class, and lie outside the scope of
$(\mathbf{O_2})$ as stated. See Lemma~\ref{lem:chi-top-bcov-split}.

\emph{Sasaki--Einstein link $\zeta$-regularisation on non-compact
toric CY$_3$.} On a non-compact CY$_3$ $X_{\mathrm{con}}$ whose
Ricci-flat Kähler cone metric approaches a Sasaki--Einstein link
$L^5$ at infinity (Klebanov--Witten
\texttt{arXiv:hep-th/9807080}: the resolved conifold $\bY$ has
$L^5=T^{1,1}\simeq S^2\times S^3$; a generic non-compact toric
CY$_3$ has $L^5$ a five-dimensional Sasaki--Einstein manifold), the
bulk $H^1(X_{\mathrm{con}},\cO)=0$ forces the classical
Atiyah--dilaton cocycle to zero in the bulk. The one-loop BCOV
diagram consequently sees only the \emph{boundary-link}
contribution, regularised by a $\zeta$-function on the link
eigenvalues of the holomorphic Laplacian. The boundary integrand
is a fermionic $\beta\gamma$-ghost system at central charge $c=-2$
on the transverse $3$-dimensional direction (Kausch
\texttt{arXiv:hep-th/9510149} Eq.~(4.12): the weight-$(1,0)$
$\beta\gamma$-system has $c=-2$), and yields a finite residual
curving $a_{\mathrm{BV}}(X_{\mathrm{con}})$ obtained by
$\zeta(0)$-regularisation of the link spectrum. For the conifold
the ghost-mode residue is $a_{\mathrm{BV}}(\bY)=-1$; this
is the exact value required to complete the Polyakov ghost-mode
balance of Remark~\ref{rem:kappa-ch-bv-distinction}.

\emph{Logical independence from the toric-rigidity axis.}
The invariant $\chi(X)/24$ is \emph{topological}: determined by the
Chern classes $c_1(T_X) = 0, c_2(T_X), c_3(T_X)$ alone, via
$\chi(X) = \int_X c_3(T_X)$. A topological invariant of a smooth
CY$_3$ is read from the de~Rham cohomology ring, an object
insensitive to the existence of a torus action, a tilting bundle, or
a toric atlas on $X$. Conversely, the obstructions $(\mathbf{O_1}),
(\mathbf{O_3}), (\mathbf{O_4})$ are statements about
automorphism-group structure ($\Aut^0(X) \supseteq T$) and
derived-categorical structure (existence of a finite-Jacobi tilting
object); both are \emph{geometric/categorical} data, not
topological data. Concretely: deformation-equivalent smooth projective
CY$_3$ share $(c_2, c_3)$ and hence $\chi$, so share the
$(\mathbf{O_2})$ value, yet can differ on Picard rank, $\Aut^0$-type,
and derived-Morita class (jumping automorphism loci inside moduli);
symmetrically, $\chi$ varies across CY$_3$ with identical
automorphism signature. Hence $(\mathbf{O_2})$ is logically
independent of the toric-rigidity axis $(\mathbf{O_1}) \Leftrightarrow
(\mathbf{O_3}) \Rightarrow (\mathbf{O_4})$ established below: neither
implies the other, and the BCOV class and the toric-rigidity data
cannot be derived from one another by any natural construction.
\end{proof}

\begin{remark}[Anomaly vanishing on retractable local models]
\label{rem:anomaly-vanishing}
A non-compact CY$_3$ $X$ admitting a proper holomorphic deformation
retraction $X \xrightarrow{r} Z$ onto a closed complex submanifold
$Z \subset X$ of complex dimension $\leq 1$ has
$\alpha_{\mathrm{BCOV}}(X) = 0$ as a class in $H^{0,1}(X, \cO_X)$.
\emph{Proof:} the retraction induces a quasi-isomorphism
$r^\ast: H^{0,\bullet}(Z, \cO_Z) \xrightarrow{\sim}
H^{0,\bullet}(X, \cO_X)$ (deformation invariance of Dolbeault
cohomology; Costello--Gwilliam vol.~I Appendix~A), and any smooth
complex curve $Z$ has $H^{0,1}(Z, \cO_Z) = 0$ exactly when $Z$ is
simply connected (e.g., $Z = \bP^1$; the conifold case) or is a
disjoint union of rational curves (toric fan $1$-skeletons without
compact $2$-cycles). When $Z$ is a positive-genus curve,
$H^{0,1}(Z, \cO_Z) \neq 0$ but the Costello--Li class still vanishes
because the Atiyah-class input $[\Omega_X]^{0,1}$ has no component
along the genus direction; the CY trivialisation of $\omega_X$
pulls $[\Omega_X]^{0,1}$ to zero on any holomorphically-embedded
curve whose normal bundle has vanishing determinant. On toric CY$_3$
without compact $4$-cycles the retract $Z$ is a union of $\bP^1$'s,
so the first mechanism suffices. Primary source: Costello--Li
\texttt{arXiv:1606.00365} Prop.~5.2 and the retraction arguments in
Costello \texttt{arXiv:1605.09656} \S 12.
\end{remark}

\begin{lemma}[Split of topological and cohomological factors]
\label{lem:chi-top-bcov-split}
The class $\alpha_{\mathrm{BCOV}}(X) = (\chi_{\mathrm{top}}(X)/24)
\cdot [\Omega_X]^{0,1} \in H^{0,1}(X, \cO_X)$ factors as a
topological scalar times a cohomology class, and each factor may
vanish independently:
\begin{enumerate}[leftmargin=2em]
\item \emph{(Resolved conifold $\bY$)} $\chi_{\mathrm{top}}(\bY) = 2
 \neq 0$ but $[\Omega_\bY]^{0,1} = 0$ (retraction; Remark
 \ref{rem:anomaly-vanishing}); $\alpha_{\mathrm{BCOV}}(\bY) = 0$.
\item \emph{(Strict compact CY$_3$, e.g., quintic)}
 $h^{0,1}(X) = 0$ by strict-CY Hodge so $[\Omega_X]^{0,1}$ lives in
 the trivial group, giving $\alpha_{\mathrm{BCOV}}(X) = 0$
 trivially, yet $\chi_{\mathrm{top}}(X) \neq 0$ generically (quintic:
 $\chi(X_5) = -200$). The relevant obstruction on strict CY$_3$ is
 the higher-order BCOV class, not the leading one.
\item \emph{($K3 \times E$)} $\chi(K3 \times E) = 0$ (Künneth on
 $\chi(E) = 0$), yet $h^{0,1}(K3 \times E) = h^{0,1}(E) = 1$ so
 $[\Omega_{K3 \times E}]^{0,1} \neq 0$. The product of zero and
 non-zero is zero: $\alpha_{\mathrm{BCOV}}(K3 \times E) = 0$ for
 dimensional reasons, not for cohomological reasons.
\item \emph{(Generic compact CY$_3$ outside these three classes)}
 requires $h^{0,1}(X) \neq 0$ (i.e., $X$ contains a
 positive-dimensional abelian factor in Bogomolov--Beauville) to
 carry a non-trivial leading anomaly.
\end{enumerate}
The sharpness of $(\mathbf{O_2})$ as a genuine compact-CY$_3$
obstruction is therefore restricted to compact CY$_3$ with a
positive-dimensional abelian factor and non-zero $\chi$; for strict
CY$_3$ the obstruction is of higher weight.
\end{lemma}

\begin{remark}[$a_{\mathrm{BV}}$ as BV anomaly coefficient; Polyakov
ghost-mode balance on class-G free-field CY$_3$]
\label{rem:kappa-ch-bv-distinction}
The BCOV one-loop curving generates a BV anomaly coefficient,
\[
 a_{\mathrm{BV}}(X)
 \;=\; \zeta(0)\big|_{L^5(X)} \cdot \mathrm{Res}_{\bar\partial}\,
 \alpha_{\mathrm{BCOV}}(X),
\]
the $\zeta(0)$-regularised boundary residue of the Costello--Li
Atiyah--dilaton class along the Sasaki--Einstein link $L^5(X)$ of a
non-compact CY$_3$, or the BV counter-term obstruction on a compact
CY$_3$. This coefficient is \emph{distinct} from the four invariants of the
subscript discipline:
\begin{itemize}[leftmargin=2em]
 \item $\kappa_{\mathrm{ch}}(X) = \sum_q (-1)^q h^{0,q}(X)$
 (chiral-supertrace, Hodge-theoretic);
 \item $\kappa_{\mathrm{cat}}(X) = \chi(\cO_X)$
 (categorical Euler, Künneth-multiplicative);
 \item $\kappa_{\mathrm{BKM}}(\Phi_N) = c_N(0)/2$
 (Borcherds/BKM weight);
 \item $\kappa_{\mathrm{fiber}}(X)$
 (lattice/fibre correction);
\end{itemize}
The BV anomaly coefficient $a_{\mathrm{BV}}$ is a separate functional
tracking the boundary ghost-mode residue of the BCOV cocycle. It is
not a fifth $\kappa_\bullet$-invariant. Unsubscripted $\kappa$ is
forbidden throughout.

\emph{Polyakov ghost-mode balance on class-G free-field CY$_3$ with
$\chi_{\mathrm{top}}=2$.} On a Polyakov free-field CY$_3$ (the
class-G cone-geometry family: resolved conifold $\bY$ is the
basic local example, $\chi_{\mathrm{top}}(\bY)=2$,
$L^5(\bY)=T^{1,1}$, Polyakov \emph{Gauge Fields and Strings} Ch.~9
free-field fermion normal ordering on the transverse
$\beta\gamma$-ghost) the local DT scalar and the BV coefficient
satisfy the \emph{renormalised ghost-mode balance law}
\[
 \kappa_{\mathrm{ch}}^{\mathrm{loc}}(\bY) + a_{\mathrm{BV}}(\bY)
 \;=\; 0.
\]
On the resolved conifold this reads
\[
 \bigl(\kappa_{\mathrm{ch}}^{\mathrm{loc}},a_{\mathrm{BV}}\bigr)(\bY)
 \;=\; (1,-1).
\]
Szendr\H{o}i two-chart DT supertrace gives
$\kappa_{\mathrm{ch}}^{\mathrm{loc}}(\bY)=1$, and the
$\zeta(0)$-regularised $T^{1,1}$ ghost-mode residue at $c=-2$
gives $a_{\mathrm{BV}}(\bY)=-1$. There is no compact
$\kappa_{\mathrm{cat}}(\bY)=\chi(\cO_{\bY})$ identity: $\bY$ is
noncompact and $\Gamma(\bY,\cO_{\bY})$ contains all fibre polynomials.
There is also no conifold BKM scalar; $\kappa_{\mathrm{BKM}}$ is defined
only after a named Borcherds automorphic input is supplied. The balance holds
\emph{only} for class-G free-field CY$_3$ at
$\chi_{\mathrm{top}}=2$ with a single compact curve class; it does
\emph{not} extend to the Theorem~C derived-centre ceiling
$K^{\kappa_{\mathrm{ch}}}\in\{0,8,13,250/3,25/3\}$, which is a separate statement
about the five-archetype
$\mathsf{G}/\mathsf{L}/\mathsf{C}/\mathsf{M}/\mathsf{B}$ landmark,
nor to generic toric CY$_3$ where the three terms decouple by the
boundary link geometry. The conflation of
$a_{\mathrm{BV}}$ with $\kappa_{\mathrm{ch}}$, and of the
conifold balance with the Theorem~C ceiling, is precisely what
the four-$\kappa_\bullet$ discipline plus a separately named BV coefficient
prevents.
\end{remark}

\begin{proof}[Proof of $(\mathbf{O_3})$]
We require an algebraic torus $T \curvearrowright X$ with isolated
fixed points on each Hilbert scheme $\cM^n(X) = \Hilb^n(X)$. A necessary
condition is that $\Aut^0(X) \supseteq T$ contain an algebraic torus.

\emph{Bogomolov--Beauville decomposition.}
Bogomolov (\emph{Hamiltonian Kähler manifolds}, Nauchn.\ Tr.\ Mat.\
1974) and Beauville (\emph{Variétés Kählériennes dont la première
classe de Chern est nulle}, J.\ Differential Geom.\ 18 (1983), 755--782,
\S2 Thm.~1) prove that any compact Kähler manifold $X$ with $c_1(X) =
0$ admits a finite étale cover splitting isometrically as
\[
 \widetilde X \;\simeq\; T_\mathrm{ab} \times \prod_i X_i^{\mathrm{CY}}
 \times \prod_j X_j^{\mathrm{IHS}}
\]
with $T_\mathrm{ab}$ a complex torus (abelian variety in the projective
case), each $X_i^{\mathrm{CY}}$ a strict Calabi--Yau ($h^{k,0}(X_i)
= \delta_{k,0} + \delta_{k,\dim X_i}$, holonomy
$\mathrm{SU}(\dim X_i)$), and each $X_j^{\mathrm{IHS}}$ an
irreducible hyperkähler (holonomy $\mathrm{Sp}(n)$). Applied to
$\dim X = 3$, the hyperkähler factor is absent (ruled out by dimension,
which must be even), and the decomposition reduces to
$\widetilde X \simeq T_\mathrm{ab} \times \prod_i X_i^{\mathrm{CY}}$
with $\dim T_\mathrm{ab} + \sum\dim X_i = 3$. The abelian summand is
detected by $h^{1,0}$: Hodge-additivity on the finite étale cover
yields $h^{1,0}(\widetilde X) = \dim T_\mathrm{ab} + \sum_i
h^{1,0}(X_i^{\mathrm{CY}}) = \dim T_\mathrm{ab}$, because strict CY
factors satisfy $h^{1,0} = 0$ by the Hodge-diamond condition. Hence
the three geometric cases partition by $h^{1,0}(X) \in \{0, 1, 2, 3\}$:
\begin{itemize}[leftmargin=2em]
\item $h^{1,0}(X) = 0$: strict CY$_3$ (no abelian factor), $\Aut^0(X)$
 finite, the Matsumura-Bogomolov-Beauville cell in which the quintic,
 bicubic, Schoen reside;
\item $h^{1,0}(X) = 1$: a mixed product such as $K3 \times E$ (strict
 CY$_2$-factor $K3$ times $1$-dim abelian factor $E$);
\item $h^{1,0}(X) \geq 2$: higher-dim abelian factor, degenerating at
 $h^{1,0} = 3$ to the pure abelian threefold.
\end{itemize}
Only the $h^{1,0}(X) = 0$ case is a \emph{strict} CY$_3$; this is the
case relevant for $(\mathbf{O_1})$--$(\mathbf{O_4})$ as stated with the
hypothesis $h^{2,0} = 0$, $\Pic^0 = 0$. The $K3 \times E$ case
(treated separately in \S\ref{sec:k3e-positioning}) sits outside the
$2.5$-count hypothesis and requires the five-fold analysis of
Proposition~\ref{prop:k3e-five-obstructions}.

\emph{Matsumura rigidity.}
Matsumura (\emph{On algebraic groups of birational transformations},
Proc.\ Japan Acad.\ 39 (1963), 181--186, Thm.~1) establishes that for
$X$ a smooth projective variety with $\omega_X$ nef, the automorphism
group scheme $\Aut(X)$ has connected component $\Aut^0(X)$ an abelian
variety: no non-trivial algebraic torus embeds in $\Aut^0(X)$. On a
CY$_3$ the Serre-triviality $\omega_X \simeq \cO_X$ supplies the
isomorphism
\[
 T_X \;\simeq\; \Omega^2_X \;\otimes\; \omega_X^\vee
 \;\simeq\; \Omega^2_X,
\]
obtained by contracting $\Lambda^2 T_X \simeq \Omega^1_X \otimes
\omega_X^\vee$ with the nowhere-vanishing section of $\omega_X$; hence
$h^0(T_X) = h^0(\Omega^2_X) = h^{2,0}(X)$ (Hodge duality), and this is
the correct Hodge identification; not $h^{1,0}$, which would be the
naive reflex. For a strict CY$_3$ ($X_i^{\mathrm{CY}}$ in the
decomposition above),
$H^0(X_i, T_{X_i}) = H^0(X_i, \Omega_{X_i}^2) = 0$ by the strict-CY
Hodge condition $h^{2,0}(X_i) = 0$, so $\Aut^0(X_i) = 1$.

\emph{Consequence for the fixed-point geometry.}
The decomposition forces one of three cases:
\begin{enumerate}[leftmargin=2em]
\item $X$ is strict CY$_3$: $\Aut^0(X)$ finite, no algebraic torus
 acts, so no equivariant localisation on $\Hilb^n(X)$.
\item $X \sim T_\mathrm{ab} \times X'$ with $\dim T_\mathrm{ab} \geq 1$:
 the acting torus $T \subseteq T_\mathrm{ab}$ has positive-dimensional
 fixed locus on $\Hilb^n(X)$ because $T_\mathrm{ab}$ translates act on
 every point of $\Hilb^n(X)$ with positive-dimensional orbits in the
 $T_\mathrm{ab}$-direction.
\item $X \sim T_\mathrm{ab}$ (abelian threefold): the entire
 $3$-dimensional torus acts, but its fixed-point set on $X$ is empty
 (translations are fixed-point-free) and hence the induced action on
 $\Hilb^n(X)$ has no isolated fixed points.
\end{enumerate}

In every case, isolated fixed points on $\Hilb^n(X)$ fail. The quintic
$X_5$ sits in case~(i) with $\Aut^0(X_5) = 1$; $K3 \times E$
(Bogomolov--Beauville-decomposed as $K3 \times E$ already; note the
Beauville factorisation separates the K3 hyperkähler from the elliptic
abelian factor, but in $\dim 3$ the K3 does not appear as a
hyperkähler summand since $\dim K3 = 2$, so $K3\times E$ decomposes as
$\dim K3 = 2$ strict CY factor times $\dim E = 1$ abelian factor) sits
in case~(ii) with $\Aut^0(K3 \times E) = E$ acting by translation,
positive-dimensional fixed locus on $\Hilb^n(K3 \times E)$. Abelian
threefold $A^3$ sits in case~(iii).

Hence Nekrasov--Okounkov equivariant-localisation shuffle presentations,
which require a torus of rank $\geq 2$ with isolated $T$-fixed points
on $\Hilb^n$, are unavailable on every compact CY$_3$.
\end{proof}

\begin{proof}[Proof of $(\mathbf{O_4})$]
A finite-quiver-with-potential presentation of the perfect derived
category requires a classical tilting generator $T \in \Perf(X)$ with
$A := \End_{D^b(X)}(T) = \bC Q / (\partial W)$ a finite-dimensional
Jacobi algebra (Van den Bergh, \texttt{arXiv:math/0211064} Def.~4.1).
The gauge torus acting on such a presentation is the \emph{outer} torus
\[
 T^{\mathrm{out}}_Q \;=\; \ker\!\left(
 \bZ^{Q_1} \xrightarrow{\;\mathbf 1\;} \bZ\right) \otimes \bC^\ast
 \;\simeq\; (\bC^\ast)^{|Q_1| - 1},
\]
rescaling each arrow $a \in Q_1$ by a character $\chi_a$ with
$\prod_a \chi_a = 1$ (the Jacobi relations $\partial_a W = 0$ force
the overall scale on $W$ to be trivial, cutting down $\bZ^{Q_1}$ by a
diagonal). This $T^{\mathrm{out}}_Q$ is distinct from the \emph{inner}
gauge torus $\prod_{v \in Q_0} \bC^\ast \subset G_v =
\prod_v \mathrm{GL}_{d_v}$ acting on representations at each vertex by
conjugation; the inner torus is a subgroup of the Morita centraliser,
the outer torus is the automorphism group of the quiver-with-potential
datum itself. For a toric CY$_3$ the global torus $T \simeq (\bC^\ast)^3$
acting on $X$ embeds into $T^{\mathrm{out}}_Q$ via the fan data
(rank-$3$ sub-torus), and this embedding is what makes chart-wise
Nekrasov--Okounkov localisation possible. On compact CY$_3$ no such
finite quiver exists to begin with; the impossibility of the outer
torus is derivative of the non-existence of $(Q, W)$.

Van den Bergh Prop.~3.2.10 establishes the converse: any NCCR $A$ of a
Gorenstein ring $R = \Gamma(X, \cO_X)$ exhibits $X$ as a small
resolution of $\mathrm{Spec}\,R$ with tilting generator reconstructing
$A$. Bocklandt (\emph{Graded Calabi--Yau algebras of dimension 3},
J.\ Pure Appl.\ Algebra 212 (2008) 14--32,
\texttt{arXiv:math/0603558} Thm.~3.1) adds the crucial refinement: for
$A$ a $3$-Calabi--Yau graded algebra, $A = J(Q, W)$ for a superpotential
$W$ on a finite quiver $Q$ if and only if the cyclic-homology pairing
is non-degenerate and the Hilbert series is compatible with the
Gorenstein parameter.

\emph{Hartogs obstruction on compact $X$.}
If $X$ is smooth projective and $A = \End(T)$ were finite-dimensional
with a tilting $T \in \Perf(X)$, the zeroth Hochschild cohomology of
$A$ would compute $R = Z(A)$:
\[
 R \;=\; \HH^0(A) \;=\; \HH^0(\Perf(X)) \;=\; H^0(X, \cO_X) \;=\; \bC
\]
(the last equality by $X$ connected projective). Hartogs's theorem then
forces $\mathrm{Spec}\,R$ to be a point, and the Van den Bergh
$\mathrm{Spec}\,R$-morphism $X \to \mathrm{Spec}\,R = \mathrm{pt}$ to
be the structure morphism. Van den Bergh's construction produces the
tilting bundle from a resolution $X \to \mathrm{Spec}\,R'$ with $R'$
an affine Gorenstein singular variety of the same dimension as $X$;
when $R' = \bC$ the resolution target is a point, i.e., $X$ would have
dimension $0$, contradicting $\dim X = 3$.

Equivalently: Bocklandt's classification (\texttt{arXiv:math/0603558}
Thm.~3.1) shows that $3$-Calabi--Yau finite Jacobi algebras exhibit
their centre $R$ as a three-dimensional Gorenstein singularity with $X$
a crepant resolution. A compact $X$ with $\Gamma(X, \cO_X) = \bC$ has
only the point as candidate base, which cannot be a three-dimensional
Gorenstein singularity.

\emph{Implications.}
Van den Bergh + Bocklandt + Hartogs-compactness thus produce the
implication $(\mathbf{O_1}) \Rightarrow (\mathbf{O_4})$: a toric
$\bC^3$-atlas would furnish a fibration $X \to \mathrm{Spec}\,\bC[x,y,z]$
in charts, patching into a small resolution structure of an affine
toric Gorenstein threefold, supplying the tilting generator. Its
absence on compact $X$ (Hartogs:\ $R = \bC$) blocks the entire
Van den Bergh--Bocklandt construction.

Beilinson's $D^b(\bP^n) \simeq D^b(B_n\text{-mod})$ for the path algebra
$B_n$ of the Beilinson quiver requires the full strong exceptional
sequence $(\cO, \cO(1), \dots, \cO(n))$, available on $\bP^n$ but not
on general compact CY$_3$. Kuznetsov (\texttt{arXiv:1404.3143}) shows
that the Kuznetsov component of a Fano threefold is generally not
equivalent to a finite quiver, further ruling out any straightforward
compact presentation.
\end{proof}

\begin{proposition}[Sumihiro biconditional:\ $(\mathbf{O_1}) \Leftrightarrow
(\mathbf{O_3})$]
\label{prop:sumihiro-biconditional}
Let $X$ be a smooth projective threefold. Then $X$ admits a toric
atlas in the sense of $(\mathbf{C_1})$ if and only if there exists an
algebraic torus $T \simeq (\bC^\ast)^3$ acting on $X$ with finitely
many fixed points and with the fixed-point-scheme decomposition
identifying the maximal cones of a fan $\Sigma \subset \bR^3$.
\end{proposition}

\begin{proof}
\emph{$(\Rightarrow)$}:\ A toric atlas $\{U_\alpha \simeq \bC^3\}$
with $\mathrm{GL}_3$-monomial transition cocycles contains an implicit
$(\bC^\ast)^3 \subset \bC^3$ in each chart. The monomial cocycles
commute with the torus actions, so the local tori glue to a global
$T \simeq (\bC^\ast)^3$-action on $X$. The fixed points are the chart
origins (assuming standard coordinates), a finite set equal to
$|\Sigma(3)|$.

\emph{$(\Leftarrow)$}:\ This direction assembles three inputs.
Sumihiro (\emph{Equivariant completion}, J.\ Math.\ Kyoto Univ.\ 14
(1974), 1--28, Thm.~1): every normal $T$-variety admits a $T$-stable
affine open cover. Luna's étale slice theorem (\emph{Slices étales},
Bull.\ Soc.\ Math.\ France, Mém.\ 33 (1973), 81--105): at each
$T$-fixed point $x \in X$ the tangent representation $T_x X$ carries a
$T$-module structure, and there is a $T$-equivariant étale neighbourhood
$U_\alpha \ni x$ together with a $T$-equivariant étale map
$U_\alpha \to T_x X$. The classification of smooth affine toric
varieties (Fulton \S2.1): a smooth $T$-stable affine variety on which
$T \simeq (\bC^\ast)^3$ acts with an isolated fixed point is
$T$-equivariantly isomorphic to $\bC^3$ with diagonal
$T$-weights $\chi_1, \chi_2, \chi_3 \in M = \mathrm{Hom}(T, \bC^\ast)$
forming a $\bZ$-basis.

Apply to the smooth projective $X$ with $T \simeq (\bC^\ast)^3$ and
finitely many fixed points. Sumihiro yields a finite $T$-stable affine
cover $\{U_\alpha\}$, with $x_\alpha \in U_\alpha$ the unique fixed
point in each piece. Each $(U_\alpha, x_\alpha)$ is a smooth affine
$T$-variety with an isolated fixed point, so by the smooth-affine-toric
classification $U_\alpha \simeq_T \bC^3$ with diagonal action. The
Luna slice supplies the explicit $T$-equivariant étale comparison to
$T_{x_\alpha} X$ and pins down the weight basis $\{\chi_1, \chi_2,
\chi_3\}$ on each chart.

Transition cocycles $U_\alpha \cap U_\beta \to U_\alpha \cap U_\beta$
must commute with the $T$-action on both sides; a $T$-equivariant
morphism between two standard toric charts is necessarily a monomial
automorphism of $\bC[x_1, x_2, x_3]_{(\text{localised})}$ with signed
integer exponents, i.e., an element of $\mathrm{GL}_3(\bZ)$ read as a
monomial element of $\mathrm{GL}_3(\bC)$ on the torus. The cocycles
glue to a complete regular fan $\Sigma \subset N_\bR = \bR^3$ with
three-dimensional cones indexing the charts and ray generators indexing
the toric divisors (Fulton \S3.4). This is the datum of a toric atlas.
\end{proof}

\begin{remark}[\emph{Two-and-a-half scope discipline}]
\label{rem:two-and-a-half}
The conventional ``four independent obstructions'' phrasing matches the
four-feature enumeration $(\mathbf{C_1})$--$(\mathbf{C_4})$ but
conflates logical content. The structural count is:
\begin{itemize}[leftmargin=2em]
\item $(\mathbf{O_1}) \Leftrightarrow (\mathbf{O_3})$ is a single
 toric-rigidity proposition (Proposition~\ref{prop:sumihiro-biconditional},
 Sumihiro 1974 Thm.~1), contributing $1$ obstruction unit;
\item $(\mathbf{O_4})$ follows from $(\mathbf{O_1})$ by Van den Bergh
 + Bocklandt + Hartogs ($R = \bC$), contributing $0.5$ additional
 categorical content;
\item $(\mathbf{O_2})$ is independent: the Costello--Li
 Atiyah--dilaton cocycle
 $\alpha_{\mathrm{BCOV}}=(\chi_{\mathrm{top}}/24)\,\tr\,\mathrm{At}(T_X)
 \in H^1(X,\cO_X)$ is a structure-sheaf class, contributing $1$
 obstruction unit.
\end{itemize}
Total:\ $2.5$ units of independent logical content. The four named
features factor through two axes: toric-rigidity and the BCOV
structure-sheaf obstruction.
\end{remark}
\subsection{Compact CY$_3$ without $K3 \times E$ fibration structure}
\label{sec:compact-no-fibration}

Among compact Calabi--Yau threefolds, the product $K3 \times E$ is
special: it inherits the $24$ elliptic-fibration fibres of the K3 factor
and the product structure with $E$, yielding twenty-four
$I_1 \times E$ nodal loci that serve as assembly stalks for the
four-stratum gluing of \S\ref{sec:gluing:k3e}. The generic compact CY$_3$ has none
of this structure.

\begin{proposition}[No fibration-assembly on generic compact CY$_3$]
\label{prop:no-fibration-assembly}
Let $X$ be a smooth projective Calabi--Yau threefold that is not
birational to a product $S \times E$ for any K3 surface $S$ and any
elliptic curve $E$, and not birational to any elliptic fibration
$\pi: X \to S$ over a surface. Then:
\begin{enumerate}[leftmargin=2em]
\item No curve-stalk decomposition $\CoHA(X) = \bigsqcup_\alpha
 \CoHA(\text{local}_\alpha)$ at isolated loci exists;
\item The Stratum-(iv) Borcherds automorphic cocycle is absent (no
 lattice-polarised period domain picks out $X$ among CY moduli);
\item Even categorical CoHA via MNOP + Davison integrality cannot be
 expressed in chart-wise form.
\end{enumerate}
\end{proposition}

The proposition is straightforward once the $K3 \times E$ mechanism is
understood. The K3 $\to \bP^1$ elliptic fibration
$\pi: K3 \to \bP^1$ with generic fibre an elliptic curve has $24$
singular fibres of Kodaira type $I_1$ (nodal rational curves) by
$\chi(K3) = 24 = \sum_{I_1} \chi(I_1)$. Multiplying by $E$ produces
$24$ loci of the form $I_1 \times E$ inside $K3 \times E$. Each such
locus is a curve-supported degenerate threefold: not a smooth $\bC^3$,
but a stalk of the form $\{xy = 0\} \times E$ embedded in a smooth
ambient $\bC^3$-neighbourhood times $E$.

For a compact CY$_3$ with no elliptic fibration; the quintic $X_5
\subset \bP^4$ being the paradigmatic example; there are no
curve-stalks to assemble. The moduli of BPS objects on $X_5$ is controlled
globally, without preferred local pieces.

\begin{example}[Quintic $X_5 \subset \bP^4$]
\label{ex:quintic}
The quintic threefold $X_5 = \{f_5 = 0\} \subset \bP^4$ has Hodge numbers
$h^{1,1} = 1$, $h^{2,1} = 101$, Euler characteristic $\chi = 2(1 - 101)
= -200$. It admits no fibration by surfaces (the Picard rank is $1$, so
any fibration base would also have Picard rank $\leq 1$; the generic
fibre of a non-trivial morphism to a curve or surface would have Picard
rank $0$, impossible for a smooth projective variety). Hence no
$24$-stalk assembly and no Stratum-(iv) Borcherds lift.

The Gromov--Witten / Donaldson--Thomas generating functions of $X_5$
are known (mirror symmetry computations of Candelas--de la Ossa--Green--Parkes;
categorical DT of Maulik--Pandharipande--Thomas), but they are
formulated globally: generating functions of degree-$d$ maps to $\bP^4$
with $f_5 = 0$, not a chart-by-chart sum. There is no local presentation.
\end{example}

\begin{example}[Bicubic $X_{3,3} \subset \bP^2 \times \bP^2$]
\label{ex:bicubic}
The bidegree $(3,3)$ hypersurface $X_{3,3} \subset \bP^2 \times \bP^2$
has $h^{1,1} = 2$ and $h^{2,1} = 83$. The projection to each $\bP^2$
factor has generic fibre a cubic curve in $\bP^2$ (elliptic curve),
giving a genus-$1$ fibration over a surface. What is absent is the
$K3\times E$ product mechanism: there is no K3 elliptic surface whose
twenty-four $I_1$ fibres become $I_1\times E$ curve-stalks after
product with an elliptic factor, and no Mukai-lattice
$\mathrm{II}_{4,20}$ period domain carrying the Gritsenko
$\Delta_5$ cocycle.

The bicubic thus sits between the quintic (no fibration structure) and
$K3 \times E$ (full elliptic fibration times $E$). It inherits a
partial Stratum-(iii) structure from the $S_3 \times S_3$ symmetry of
the $\bP^2 \times \bP^2$ coordinates, and partial Stratum-(iv) from
the bicubic period domain, but no chart-wise $Y^+(\widehat{\fgl}_1)$
assembly.
\end{example}

\begin{example}[Schoen threefold]
\label{ex:schoen}
The Schoen threefold is a smooth fibre product of two rational elliptic
surfaces over $\bP^1$: $\widetilde X = S_1 \times_{\bP^1} S_2$ with
$S_i \to \bP^1$ extremal rational elliptic surfaces. It has
$h^{1,1} = 19$, $h^{2,1} = 19$, $\chi = 0$. The Schoen threefold
admits \emph{two} elliptic fibrations (one per factor), giving richer
Stratum-(ii)/(iv) structure than the quintic but less than
$K3 \times E$. Stratum-(iv) automorphic gluing via the Narain lattice
$\mathrm{II}_{1,17} \oplus \mathrm{II}_{1,17}$ is theoretically available;
chart-wise formulas remain out of reach.
\end{example}

The pattern of Propositions~\ref{prop:no-fibration-assembly} across the
census of compact CY$_3$ is simple: the \emph{less} fibration structure,
the fewer strata the geometry occupies, the fewer gluing cocycles are
available. The quintic is the minimal case (only MNOP categorical CoHA,
no chart-wise pieces); $K3 \times E$ is the maximal case (four strata
intersect, four cocycles assemble). Partial-fibration examples such as
the bicubic and Schoen threefold lie between these endpoints; a strict
generic compact CY$_3$ lies at the minimal end.
\subsection{Partial transfer on compact CY$_3$}
\label{sec:partial-transfer}

The chart-wise construction of Parts~II--VII terminates at the
toric boundary, while several stratum-free receptacles remain. They are
targets and conditional packages, not a compact-CY$_3$ closure theorem.
Their exact hypotheses govern compact-CY$_3$-native CoHA constructions.

\begin{proposition}[Stratum-free receptacles and obligations]
\label{prop:stratum-free}
Let $X$ be any smooth projective CY$_3$ (compact, not assumed toric or
fibred). The following structures are available at the indicated scope:
\begin{enumerate}[leftmargin=2em]
\item The \emph{Davison--Meinhardt critical-Hall receptacle}
 $\CoHA_{\mathrm{crit}}(X) := H^\bullet_{\mathrm{crit}}(\cM_X^{ss},
 \phi_{W})$ once the compact moduli problem is supplied with local
 $d$-critical charts, orientation data, and proper Hall
 correspondences. This is a stratum-free hypotheses, not a
 consequence of compactness alone; primary source
 Davison--Meinhardt \texttt{arXiv:1601.02479}.
\item The \emph{Kontsevich--Soibelman motivic wall-crossing formula}
 in the motivic Hall algebra, with its quantum-dilogarithm cocycle and
 chamber-to-chamber automorphism, after a stability condition,
 orientation data, and integration map have been fixed; primary source
 \texttt{arXiv:0811.2435}.
\item The \emph{two-stage framed object-level problem}
 $\Phi_3^{(\Sigma_2,C)} =
 \mathrm{Sp}^{\mathrm{ch}}_{\Sigma_2, C} \circ \Phi^{FA}_3$.
 Stage~1 is a principle-level factorisation target after the
 Kontsevich--Tamarkin formality/associator torsor point and
 Costello--Gwilliam--Li locality are chosen. Stage~2 is a framed
 specialisation problem requiring a witnessed admissible
 $(\Sigma_2,C)$ datum, a chain-level $\bS^3$-framing witness, and the
 analytic-completion inputs of Theorem~\ref{thm:cy-to-chiral-d3}.
 Dunn--Lurie additivity identifies the operadic target; it does not
 supply these witnesses.
\end{enumerate}
\end{proposition}

\begin{proof}[Sketch]
Each item is independent of a toric atlas, torus localisation, and a
finite quiver with potential. Independence from toric geometry is not
automatic existence on every compact $X$: each item carries its own
orientation, realisation, convergence, or framing hypotheses.

Under the Davison--Meinhardt hypotheses, critical integrality produces
a BPS Lie algebra $\fg_X^{\mathrm{BPS}}$ graded by
$K^{\mathrm{num}}(X)$ and a PBW-type vector-space factorisation of the
critical Hall algebra. This is not a shuffle presentation, not a finite
quiver presentation, and not a no-extra-relations theorem for a
Borcherds or Hall--Drinfeld double.

Motivic wall-crossing (item 2) is formulated in the motivic Hall
algebra $\mathrm{MH}(\cM_X^{ss})$, with quantum-dilogarithm variables
$\Psi(x_\gamma)$ indexed by numerical classes. The formula
$\prod_{\mathrm{incoming}} \Psi(x_\gamma)^{\Omega(\gamma)} =
\prod_{\mathrm{outgoing}} \Psi(x_\gamma)^{\Omega(\gamma)}$ holds
chamber-by-chamber once the stability and motivic-integration data have
been fixed; it supplies numerical wall-crossing, not a compact
Hall--Drinfeld presentation.

Two-stage factorisation (item 3) is a categorical target. The
principle-level Stage~1 invokes Kontsevich--Tamarkin $E_3$-formality of
$\HH^\bullet(\Coh\,X)$ plus Costello--Gwilliam--Li holomorphic locality.
The framed output $\Phi_3^{(\Sigma_2,C)}$ exists only after the
witnessed $d=3$ hypotheses are supplied. Therefore the proposition
does not assert a compact $\Phi_3$ for an arbitrary quintic or generic
compact CY$_3$.
\end{proof}

\begin{remark}[Framing diagnostics are not compact closure]
\label{rem:framing-diagnostics-not-closure}
Dunn--Lurie additivity, the Goodwillie filtration, topological
$\bS^3$-framing vanishing, \v{C}ech--HTT convergence, and Borel
discriminants identify obstruction targets and sufficient criteria.
They do not imply universal
$\HH^{-2}_{E_1}(A,A)=0$, contractible $\bS^3$-lifting spaces, compact
$\Phi_3$, realised compact Hall/CoHA, PBW presentations, or
no-extra-relations.

The strict chain-level obstruction already separates the raw and
corrected operators:
\[
 [m_3,B_{\mathrm{term}}^{(2)}][a|a|a|a|b]
 =
 2\alpha[b]\neq 0 .
\]
Thus raw termwise cancellation for $B_{\mathrm{term}}^{(2)}$ is false.
The raw operator $B_{\mathrm{term}}^{(2)}$ is not Costello's corrected
operator $B_{\mathrm{TCFT}}^{(2)}$ unless a strictifying homotopy and a
comparison datum have been supplied. Costello's corrected operator
belongs to a different TCFT comparison. Closure requires either a
comparison from the corrected TCFT datum to the $\bS^3$-framing
obstruction complex, or a complete filtered theorem for
$\HH^{-2}_{E_1}$: a named comparison map, a complete, exhaustive,
separated, strongly convergent filtration, and an empty total-degree
$-2$ line.
\end{remark}

\begin{proposition}[What fails to transfer]
\label{prop:transfer-failures}
The following chart-wise machinery does not transfer to a generic
compact CY$_3$:
\begin{enumerate}[leftmargin=2em]
\item Per-chart Schiffmann--Vasserot $Y^+(\widehat{\fgl}_1)$ stalks
 located at toric vertices;
\item Galakhov--Li--Yamazaki bond factors $\varphi_{a \Rightarrow b}$
 attached to toric fan edges;
\item The shuffle-algebra presentation of $\CoHA(X)$ via
 Nekrasov--Okounkov fixed-point localisation;
\item The Van den Bergh tilting bundle producing a finite quiver with
 potential.
\end{enumerate}
\end{proposition}

Each item requires at least one of $(\mathbf{C_1})$--$(\mathbf{C_4})$,
and Theorem~\ref{thm:four-obstructions} rules them out. The Stage-1
factorisation target $\Phi^{FA}_3$ can still be formulated, but lacks
the explicit local formula that toric geometry provides; such
principle-level formulation does not yield a computation.

\begin{remark}[The meaning of \emph{principle-level}]
\label{rem:principle-level}
Proposition~\ref{prop:stratum-free} item (3) is ``principle-level'': the
Stage~1 $\Phi^{FA}_3$ is an abstract target after
Kontsevich--Tamarkin plus Costello--Gwilliam--Li inputs are fixed, but
the framed object-level specialisation
$\Phi_3^{(\Sigma_2,C)}$ requires H1--H4 and a fixed admissible datum
$(\Sigma_2,C)$. The Stage-1 output $\Phi^{FA}_3(X)$ is not computable
by chart-wise formulas on an arbitrary compact CY$_3$.
Computing it requires either (a) a specific tilting-bundle presentation
(which $(\mathbf{O_4})$ blocks), or (b) a specific local chart
presentation (which $(\mathbf{O_1})$ blocks), or (c) a specific torus
equivariance (which $(\mathbf{O_3})$ blocks). None is available for the
quintic.

What is available for the quintic: the Davison--Meinhardt
critical-Hall receptacle is the correct categorical target once the
compact orientation, realisation, support, and properness package is
supplied; its \emph{generating function} of BPS numbers is computed by
mirror symmetry; the link between categorical classes and numerical
invariants is the Joyce--Song integration map
(\texttt{arXiv:0810.5645}). The algebra structure of
$\CoHA_{\mathrm{crit}}(X_5)$ remains non-explicit.
\end{remark}
\subsection{What is recovered on compact CY$_3$}
\label{sec:recovery}

The collapse of chart-wise gluing on compact CY$_3$ is not the end of
the story. Three non-chart-wise structures remain at their proper
scope: conditional BPS Lie algebra multiplicity data, the MMNS motivic
generating function, and the Borcherds--Gritsenko automorphic formula
on lattice-polarised strata. They recover numerical or target-side
content without producing an explicit compact cocycle.

\begin{proposition}[Davison BPS Lie algebra]
\label{prop:davison-bps}
\emph{(Davison \texttt{arXiv:1601.02479} Thm.~A;
Davison--Meinhardt \texttt{arXiv:1601.02479})}
For a CY$_3$ critical-Hall moduli problem attached to $X$, after the
orientation, local critical-chart, realisation, and Hall-correspondence
hypotheses have been supplied, there exists a BPS Lie algebra
\[
 \fg_X^{\mathrm{BPS}} \;\subset\; \CoHA_{\mathrm{crit}}(X)
\]
graded by $K^{\mathrm{num}}(X)$, such that:
\begin{enumerate}[leftmargin=2em]
\item The PBW theorem holds at the integrality/vector-space level:
 $\CoHA_{\mathrm{crit}}(X) \simeq
 U(\fg_X^{\mathrm{BPS}})$ as graded vector space;
\item The multiplicity of $\fg_X^{\mathrm{BPS}}$ in each numerical class
 $\gamma$ is the BPS invariant $\Omega(\gamma)$.
\end{enumerate}
This is not a finite presentation of $\CoHA_{\mathrm{crit}}(X)$, not a
Drinfeld-double construction, and not a no-extra-relations theorem.
\end{proposition}

On a strict CY$_3$ the numerical Grothendieck group has rank
$2 + 2\,h^{1,1}(X)$: the Chern character identifies
$K^{\mathrm{num}}(X)_\Q$ with the algebraic part of
$H^{\mathrm{even}}(X, \Q) = H^0 \oplus H^2 \oplus H^4 \oplus H^6$,
the Euler form is non-degenerate by construction, and the strict
condition $h^{2,0} = 0$ collapses $H^2$ to the Hodge class $H^{1,1}$
of rank $h^{1,1}$, with $H^4 \simeq H^{2,2}$ dually of rank
$h^{2,2} = h^{1,1}$ (Hodge symmetry and Serre duality) and
$H^0, H^6$ each of rank $1$. For the quintic ($h^{1,1} = 1$):
\[
 K^{\mathrm{num}}(X_5) \;\simeq\; \bZ^{1 + 1 + 1 + 1} \;=\; \bZ^4,
\]
with the four factors the Chern-character components
$(\mathrm{rk}, c_1, c_2, c_3)$. The BPS invariants
$\Omega(\gamma)$ are computed by the Candelas--de la Ossa--Green--Parkes
$B$-model calculation; the BPS Lie algebra structure is encoded in the
generating function but not yet tabulated vertex-by-vertex (no finite
quiver indexes the vertices).

\begin{proposition}[MMNS motivic generating function]
\label{prop:mmns-motivic}
\emph{(Morrison--Mozgovoy--Nagao--Szendrői
\texttt{arXiv:1107.5017})}
The motivic DT partition function of a CY$_3$ category can be organised
into an infinite product over BPS classes:
\[
 Z^{\mathrm{DT,mot}}(X) \;=\;
 \prod_{\gamma} \prod_{k \geq 0}
 \left(1 - \mathbb{L}^{k - d_\gamma/2} y^\gamma\right)^{-\Omega(\gamma)
 \,\chi(L_\gamma)}
\]
for suitable line bundles $L_\gamma$ and motivic weights.
\end{proposition}

The MMNS formula expresses the DT generating function as a motivic
Euler product, with the BPS numbers $\Omega(\gamma)$ and motivic
weights $\chi(L_\gamma)$ as numerical inputs. For the conifold the
formula reduces to $Z^{\mathrm{DT}}(\bY) = M(-q)^2 \prod_{k \geq 1}(1 -
Q(-q)^k)^k$ (Szendrői's two-chart form). For compact CY$_3$ the formula
is formally available but the $\Omega(\gamma)$ values are not
closed-form.

\begin{proposition}[Borcherds--Gritsenko automorphic assembly]
\label{prop:borcherds-assembly}
\emph{(Borcherds \texttt{arXiv:alg-geom/9609022},
Gritsenko \texttt{arXiv:math/9906190})}
For a lattice-polarised Calabi--Yau threefold $X$ whose moduli sits
inside a locally-symmetric period domain $\cD / \Gamma$, there exists
a Borcherds lift $\Psi_X$ of a weak Jacobi form $\phi$ on the
polarising lattice, with
\[
 \kappa_{\mathrm{BKM}}(\Psi_X) \;=\; c_{\Psi_X}(0)/2,
\]
where $c_{\Psi_X}(0)$ is the constant Fourier coefficient of the
generating Jacobi form.
\end{proposition}

On $K3 \times E$ the lattice-polarisation is the Mukai lattice
$\widetilde H(K3, \bZ) \simeq \mathrm{II}_{4,20}$ and the Borcherds
lift is the Gritsenko $\Delta_5$ form of weight $5$ with constant term
$c_1(0) = 10$, giving $\kappa_{\mathrm{BKM}}(\Delta_5) = 5$. For the quintic
there is no period domain on which the classical Borcherds construction
produces a canonical lift. Hence the automorphic assembly is
stratum-(iv) specific; compact CY$_3$ outside stratum (iv) do not inherit
it.

\begin{theorem}[What recovers, what does not]
\label{thm:recovery-summary}
On a generic compact CY$_3$ (Picard rank one, no fibration), the
stratum-free machinery supplies only the following conditional
receptacles and numerical packages:
\begin{itemize}[leftmargin=2em]
\item $\CoHA_{\mathrm{crit}}(X)$, once the compact critical-Hall
 orientation, realisation, support, and properness package is supplied;
\item $\fg_X^{\mathrm{BPS}} \subset \CoHA_{\mathrm{crit}}(X)$ and its
 PBW vector-space integrality, under the same Davison--Meinhardt
 hypotheses, graded by $K^{\mathrm{num}}$;
\item $Z^{\mathrm{DT,mot}}(X)$ (MMNS motivic generating function);
\item $\Omega(\gamma)$ via mirror symmetry (numerical BPS invariants);
\item $\Phi^{FA}_3(X)$ as a principle-level Stage~1 target, not a
 constructed compact framed $\Phi_3^{(\Sigma_2,C)}(X)$.
\end{itemize}
Not produced:
\begin{itemize}[leftmargin=2em]
\item Chart-wise $Y^+(\widehat{\fgl}_1)$ stalks;
\item Explicit bond-factor formulas per edge of a toric fan;
\item Shuffle-algebra presentation with torus-equivariant localisation;
\item Finite quiver with potential $(Q, W)$;
\item Compact framed $\Phi_3^{(\Sigma_2,C)}(X)$ without the witnessed
 $S^3$-framing, TCFT comparison, analytic completion, and admissible
 specialisation datum;
\item Hall--Drinfeld double data, PBW presentation, and
 no-extra-relations beyond the Davison--Meinhardt vector-space
 integrality statement;
\item Borcherds--Gritsenko automorphic cocycle (absent outside stratum
 (iv)).
\item A global quantum vertex chiral group $G(X)$; constructing it
 requires the representability package of
 Chapter~\ref{ch:cy-c-six-routes-convergence},
 Definition~\ref{def:cy-c-gx-representability-datum}.
\end{itemize}
In particular, the Stage-1 target and critical-Hall receptacle are not
closed-form compact constructions, and the BPS Lie algebra, when its
hypotheses are met, is not combinatorially presented.
\end{theorem}

\begin{theorem}[No stratum-free construction of \texorpdfstring{$G(X)$}{G(X)}]
\label{thm:no-stratum-free-gx}
Let $X$ be a smooth projective compact CY$_3$. The stratum-free data
that can be formulated for such $X$, and constructed only after their
own hypotheses are supplied,
\[
 \CoHA_{\mathrm{crit}}(X),\qquad
 \fg_X^{\mathrm{BPS}}\subset \CoHA_{\mathrm{crit}}(X),\qquad
 \PhiFA_3(\Perf X),
\]
does not determine a quantum vertex chiral group $G(X)$. The missing
data are exactly the non-stratum-free entries:
\begin{enumerate}[label=\textup{(\alph*)}, leftmargin=2.5em]
\item a finite route diagram of independently constructed
$E_1$-chiral outputs and comparison maps;
\item a positive half carrying a compatible coproduct;
\item a completed negative half and a non-degenerate Hopf pairing;
\item a completion in which the Drinfeld double exists;
\item continuity of the braided centre through the route colimit.
\end{enumerate}
Thus even when a compact CY$_3$ has the critical-Hall package and a
Stage-$1$ factorisation target, the symbol $G(X)$ remains
unconstructed.
For a strict Picard-rank-one compact CY$_3$ with no fibration or
lattice-polarised period stratum, items \textup{(a)}--\textup{(e)} are
not produced by the machinery of this chapter. For $K3\times E$, items
\textup{(a)} and part of \textup{(b)} are supplied by the stratified
six-route locus; items \textup{(c)}--\textup{(e)} are the conditional
Hall--Drinfeld assembly obligations of CY-C.
\end{theorem}

\begin{proof}
Davison--Meinhardt critical cohomology supplies an associative Hall
product and a BPS Lie algebra only after its critical-Hall hypotheses
are supplied. The two-stage construction supplies an
$E_3$-holomorphic factorisation target and, after a fixed witnessed
framed specialisation, an $E_1$-chiral algebra. Neither construction
supplies a Hopf copairing, a negative half, or a completed braided
monoidal category in which the Drinfeld double and centre are defined. A
Drinfeld double is not a functor from arbitrary associative CoHAs to
quantum vertex groups; it is a construction on paired Hopf data in a
specified completion. Likewise the Drinfeld centre is a construction on
a monoidal representation category and is not forced by the existence
of the underlying $E_1$ algebra.

The first obstruction is therefore formal: even after the conditional
stratum-free data are supplied, they stop at
$\CoHA_{\mathrm{crit}}(X)$, $\fg_X^{\mathrm{BPS}}$, and
$\PhiFA_3(\Perf X)$. Proposition~\ref{prop:no-fibration-assembly} and
Proposition~\ref{prop:transfer-failures} identify why a strict generic
compact CY$_3$ does not acquire the missing route and Hopf data from the
chart-wise construction. The $K3\times E$ exception is not a global
construction; it is the named locus where the extra data are visible
and still conditional.
\end{proof}

\begin{remark}[\emph{Compact CY$_3$ outside $K3 \times E$}]
\label{rem:generic-compact}
For compact CY$_3$ that inherit a partial stratum structure; e.g.,
the bicubic with $\bP^2 \times \bP^2$ symmetry (partial Stratum~(iii))
or the Schoen threefold with two elliptic fibrations (partial
Stratum~(ii/iv)); the recovery is intermediate: more than the
quintic, less than $K3 \times E$. The automorphic cocycle may survive
in weakened form; the chart-wise stalks do not.

The DT/motivic integration package and the Davison--Meinhardt
critical-Hall receptacle are the irreducible minimum. When their
compact orientation, realisation, and support hypotheses are supplied,
they produce the BPS Lie algebra graded by $K^{\mathrm{num}}$. This is
the universal target, not a universally realised algebra. Everything
beyond requires at least one of the four features, or the presence of
lattice-polarisation.
\end{remark}
\subsection{The $K3 \times E$ positioning}
\label{sec:k3e-positioning}

$K3 \times E$ sits outside the hypothesis of
Theorem~\ref{thm:four-obstructions} ($h^{2,0} = 0$, $\Pic^0 = 0$): by
Künneth $h^{1,0}(K3 \times E) = h^{1,0}(K3) + h^{1,0}(E) = 0 + 1 = 1$,
and $\Pic^0(K3 \times E) = \mathrm{Pic}^0(E) = E$ is positive-dimensional.
The 2.5-count structural logic therefore does not apply directly to
$K3 \times E$; the failure of a global NCCR is governed instead by the
five-fold catalogue of Proposition~\ref{prop:k3e-five-obstructions}. The
threefold is atypical on three further counts relative to a generic
compact CY$_3$:

\begin{enumerate}[leftmargin=2em]
\item $h^{1,0}(K3 \times E) = 1$ places $\Aut^0(K3\times E) \supseteq E$
 as a positive-dimensional algebraic group (translations by $E$).
 Stratum~(ii) (reduced $\bC^\times \times \Aut(X)$) is non-trivial.
\item It admits the K3 elliptic fibration $K3 \to \bP^1$ pulled back to
 $K3 \times E \to \bP^1 \times E$, giving $24$ singular fibres of type
 $I_1 \times E$ over the $24$ discriminant points. Stratum~(i)
 (toric-like local structure) is partial.
\item Its periods sit in a $20$-parameter locally-symmetric domain
 $\mathcal{D} = \mathrm{O}(2,19) / \mathrm{O}(2) \times \mathrm{O}(19)$
 (K3 factor) times $\mathbb{H}$ (E-factor), admitting the
 Gritsenko $\Delta_5$ Borcherds lift. Stratum~(iv) is fully available.
\end{enumerate}

This triple alignment; positive-dimensional $\Aut^0$, elliptic
fibration structure, lattice-polarised period domain; is exceptional.
It is what permits the $24$-curve-stalk assembly of \S\ref{sec:gluing:k3e}. The
generic compact CY$_3$ (quintic, bicubic, Schoen) realises at most one
of the three.

\begin{proposition}[Five-fold obstruction to global NCCR on $K3\times E$]
\label{prop:k3e-five-obstructions}
The compact Calabi--Yau threefold $X = K3 \times E$ admits no global
non-commutative crepant resolution (NCCR) in the Van den Bergh sense.
The trace-splitting obstruction rules out a single tilting generator.
Four standard construction mechanisms fail for separate geometric
reasons:

\begin{enumerate}[label=\emph{(\alph*)}, leftmargin=2.5em]
\item \emph{Trace-splitting obstruction.} If a vector-bundle tilting
generator $T$ existed, the trace map
$\mathcal End(T)\to\cO_X$ would split after division by
$\operatorname{rank}T$. Thus $H^i(X,\cO_X)$ would be a direct summand
of $\Ext^i_X(T,T)$. The tilting axiom $\Ext_X^{>0}(T,T)=0$ would force
$H^{>0}(X,\cO_X)=0$, but $H^1(E,\cO_E)$ and $H^2(K3,\cO_{K3})$ survive
in the K\"unneth decomposition.

\item \emph{Derived McKay requires finite $\Aut$ fixing a point.}
The derived McKay correspondence of Bridgeland--King--Reid
(\texttt{arXiv:math/9908027}) constructs NCCRs on $[V/\Gamma]$ when
$\Gamma \subset \mathrm{SL}(V)$ is \emph{finite} and acts with a fixed
point. On $K3\times E$ the group $\Aut^0(K3\times E) = E$ is a
positive-dimensional connected algebraic group acting without fixed
points (translations), disqualifying the orbifold-McKay production of
an NCCR.

\item \emph{HPD self-duality incompatible with product polarisation.}
Homological projective duality (Kuznetsov \texttt{arXiv:math/0507292})
would construct NCCRs from self-dual Lefschetz decompositions of
$D^b(\Coh\,X)$ relative to a polarisation $L \to X$. On the product
$K3 \times E$ the natural polarisations factor as
$L = L_{K3} \boxtimes L_E$, breaking HPD self-duality: the Lefschetz
decomposition of $D^b(K3)$ relative to $L_{K3}$ is not compatible with
the $L_E$-Lefschetz decomposition of $D^b(E)$ in any way that induces
a Kuznetsov-self-dual structure on the product.

\item \emph{Mukai vanishing fails off the $K3$ factor.} Mukai
(\emph{Symplectic structure of the moduli space of sheaves on an
abelian or K3 surface}, Invent.\ Math.\ 77 (1984) 101--116) proves
$\Ext^i_{K3}(\cE, \cE) = 0$ for $i \neq 0, 2$ for simple rigid sheaves
on $K3$. On $K3 \times E$ the product sheaves $\cE \boxtimes \cF$ with
$\cF \in \Coh(E)$ admit $\Ext^i_{E}(\cF, \cF) \neq 0$ at $i = 0, 1$,
breaking the vanishing and preventing the Mukai-flop NCCR construction
from globalising.

\item \emph{CY-$3$ d-critical structure does not promote to a
tilting generator.} A global $d$-critical atlas on the moduli
$\cM^n(X)$; equivalently, a global choice of
$(-1)$-shifted symplectic structure on $\cM^n(X)$ in the sense of
Pantev--Toën--Vaquié--Vezzosi (\texttt{arXiv:1111.3209}) descending
via the Darboux theorem of Brav--Bussi--Joyce
(\texttt{arXiv:1305.6302}, Thm.~5.18) to a Zariski-local presentation
as the critical locus of a regular function; exists on any
CY$_3$ moduli stack. What BBJ Thm.~5.18 supplies is Zariski-local
$d$-critical charts $(U_\alpha, f_\alpha)$ with
$\cM^n(X)|_{U_\alpha} = \mathrm{crit}(f_\alpha)$; it does \emph{not}
supply a tilting generator gluing these charts to a global
$\End(T)$-representation. The $d$-critical data is by construction
chart-local, and on compact CY$_3$ without NCCR input the charts do
not patch to a global finite-dimensional algebra $A$ with
$D^b(\Coh\,X) \simeq D^b(A\text{-mod})$. The Behrend--Fantechi
symmetric obstruction theory (\texttt{arXiv:math/0512556}) is the
virtual-fundamental-class input on $\cM^n(X)$, a separate structure
from the tilting-generator question. On $K3 \times E$ the PTVV form
exists globally (pulled back from either factor) but does not lift to
a tilting object: the local-to-global passage from $d$-critical chart
data to tilting generator is the obstruction, not the availability of
local symplectic structure.
\end{enumerate}
\end{proposition}

\begin{remark}[Serre-equivariant quasi-NCCR substitute]
\label{rem:serre-equivariant}
The absence of a global NCCR on $K3 \times E$ is not absolute: the
trace-splitting obstruction and the four failed construction routes in
Proposition~\ref{prop:k3e-five-obstructions} leave open a weakened
form, the \emph{Serre-equivariant quasi-NCCR datum}: a locally-defined
tilting object $T_{\mathrm{loc}} \in D^b(\Coh\,U)$ on each affine open
$U \subset X$, equivariant under the Serre twist
$\mathrm{Serre} = - \otimes \omega_X[\dim X] \simeq [3]$, glued not by a
global module but by the factorisation-pushforward cocycle
\[
 \sigma_{\alpha\beta}:\;
 (j_{U_\alpha})_\ast T_{\mathrm{loc}}^\alpha \otimes_{\cO_{U_\alpha \cap
 U_\beta}} (j_{U_\beta})^\ast \;\longrightarrow\;
 (j_{U_\beta})_\ast T_{\mathrm{loc}}^\beta \otimes_{\cO_{U_\alpha\cap
 U_\beta}} (j_{U_\alpha})^\ast,
\]
where $j_{U_\alpha}$ denotes the inclusion and the cocycle condition
$\sigma_{\alpha\beta}\circ\sigma_{\beta\gamma}
= \sigma_{\alpha\gamma}$ on triple overlaps is checked in the homotopy
category $\mathrm{Ho}(D^b(U_\alpha\cap U_\beta\cap U_\gamma))$. The
resulting structure is \emph{quasi}-NCCR: it provides local tilting
charts for bounded Davison--Meinhardt critical Hall truncations and for
factorisation-homology $\Phi^{FA}_3$, but it does not provide a global
$\End(T)$-representation of $D^b(\Coh\,X)$ and does not construct the
completed compact CoHA unless the compact-passage defects
$o_{\mathrm{ML}},o_{\mathrm{real}},o_{\mathrm{cpt}}$ vanish. The
categorical CoHA of $K3\times E$ is assembled through the stratum-(iv)
Borcherds cocycle rather than through any finite-quiver-with-potential
presentation.
\end{remark}
\subsection{Finite Rees layer and compact Hall obstructions}
\label{sec:compact-oriented-critical-hall-cosheaves}

The Serre-equivariant quasi-NCCR keeps enough local tilting geometry to
name critical Hall charts. At finite height, for fixed bounds on jets,
renormalisation order, charge, and Ran arity, the
Dolbeault--Weiss--Ran Rees construction of
Proposition~\ref{prop:gluing-finite-dwr-ran-rees} gives a
Maurer--Cartan element on a finite \v{C}ech/Ran nerve. The compact
passage begins after this finite Rees layer: inverse limits, critical
realisation, compact-support functoriality, and the Hall--Drinfeld
double are not consequences of the Stokes sum.

\begin{definition}[Compact passage datum beyond the finite Rees layer]
\label{def:compact-oriented-critical-hall-datum}
Let $X=K3\times E$, fix a stability chamber $\sigma$ and a strict
central-charge sector $S$. Let $\mathfrak U$ be a DWR-good cover and
let $\mathcal I$ be the cofiltered category of finite bounds
$(N,r,L,m)$ on holomorphic jets, renormalisation order, charge, and
Ran arity. Assume a multiplicative finite DWR/Ran cyclic critical atlas
has been fixed at every object of $\mathcal I$. A
\emph{compact passage datum beyond the finite Rees layer} consists of
four pieces.
\begin{enumerate}[label=\emph{(\roman*)}, leftmargin=2.5em]
\item \emph{Mittag--Leffler completion.} The finite Rees Hall complexes
\[
 \mathcal H^{N,r,L,m}_{\mathrm{Rees}}(\tau)
\]
and the finite maps
$\Theta_{\hCS\to\Hall}^{\mathrm{Rees},\mathrm{or};N,r,L,m}$ form an
inverse system over $\mathcal I$. The completion defect is the
derived-limit class
\[
 o_{\mathrm{ML}}
 \in
 R^{1}\!\varprojlim_{\mathcal I}\,
 H^\bullet\!\left(
 \mathcal H^{N,r,L,m}_{\mathrm{Rees}}(\tau)
 \right),
\]
or equivalently the failure of local finiteness in the
Harder--Narasimhan charge/radius filtration
\[
 \widehat{\bigoplus}_{\gamma\in\Gamma_{\sigma,S}(\tau)}
 :=
 \varprojlim_R
 \bigoplus_{\substack{\gamma\in\Gamma_{\sigma,S}(\tau)\\
 |Z_\sigma(\gamma)|\leq R}}
\]
to be an exact inverse-limit operation compatible with refinements and
wall crossing.
\item \emph{Critical realisation.} A monoidal realisation functor
\[
 \mathcal R_{\mathrm{crit}}\colon
 \Hall_{\sigma}^{\mathrm{Rees},\mathrm{or}}
 \longrightarrow
 H^{BM}_{G_{\tau,\gamma}}\!
 \left(
 \operatorname{Crit}(f_{\tau,\gamma}),
 \phi_{f_{\tau,\gamma}}\otimes\mathscr L_{\tau,\gamma}
 \right)[s(\tau,\gamma)](t(\tau,\gamma))
\]
compatible with Thom--Sebastiani, lci pullback, equivariant
Borel--Moore pushforward, Joyce--Kontsevich--Soibelman orientations,
shifts, Tate twists, and the chosen completion. Its defect is denoted
$o_{\mathrm{real}}$.
\item \emph{Compact support.} The Hall extension correspondences on
semistable compactly-supported perfect objects
\[
 \mathcal M_{\tau,\gamma_1}^{ss}\times
 \mathcal M_{\tau,\gamma_2}^{ss}
 \xleftarrow{\ p_1\ }
 \mathcal E_\tau(\gamma_1,\gamma_2)
 \xrightarrow{\ p_2\ }
 \mathcal M_{\tau,\gamma_1+\gamma_2}^{ss}
\]
must satisfy the support and properness conditions required for
$(p_2)_!p_1^!$ after refinement and completion. The failure of these
conditions is $o_{\mathrm{cpt}}$.
\item \emph{Double data.} The realised positive half
$Y^+=\CoHA^{or}_{crit}(X)_{\sigma,S}^{\wedge}$ must be equipped with a
completed Serre-dual negative half $Y^-$, a Cartan completion $Y^0$, a
continuous coproduct $\Delta_+$, a non-degenerate Hopf pairing
$\langle-,-\rangle:Y^+\widehat\otimes Y^-\to\bC((\hbar))$,
cross-relations, radical quotient, and centre compatibility
$Z_{\mathrm{cent}}$. The corresponding defects are
$o_\Delta$, $o_{\mathrm{pair}}$, and $o_{\mathrm{cent}}$.
\end{enumerate}
\end{definition}

\begin{proposition}[Finite compact Hall product criterion]
\label{prop:finite-compact-hall-product-gate}
Fix a finite compact-support height window and three retained charges
\(\gamma_1,\gamma_2,\gamma_3\). Let
\[
 H_i,\quad H_{ij},\quad H_{123}
\]
be the corresponding finite Borel--Moore vanishing-cycle spaces, with
compact-support projection matrices
\[
 P_i,\quad P_{ij},\quad P_{123}.
\]
Suppose finite Hall product matrices have been supplied:
\[
 m_{12}\colon H_1\otimes H_2\to H_{12},\qquad
 m_{23}\colon H_2\otimes H_3\to H_{23},
\]
\[
 m_{12,3}\colon H_{12}\otimes H_3\to H_{123},\qquad
 m_{1,23}\colon H_1\otimes H_{23}\to H_{123}.
\]
Let \(T_L,T_R\) be the two Thom--Sebastiani transport matrices for
the two parenthesisations, and let \(O_L,O_R\) be the corresponding
orientation-line transport matrices. The supplied finite compact Hall
product is associative on the retained compact-support window exactly
when
\[
 m_{12,3}(m_{12}\otimes 1_{H_3})
 =
 m_{1,23}(1_{H_1}\otimes m_{23}),
 \qquad
 T_L=T_R,\qquad O_L=O_R,
\]
and when the compact-support projections are genuine projections and
intertwine the four products:
\[
 P_a^2=P_a\quad
 (a=1,2,3,12,23,123),
\]
\[
 P_{12}m_{12}=m_{12}(P_1\otimes P_2),\qquad
 P_{23}m_{23}=m_{23}(P_2\otimes P_3),
\]
\[
 P_{123}m_{12,3}=m_{12,3}(P_{12}\otimes P_3),\qquad
 P_{123}m_{1,23}=m_{1,23}(P_1\otimes P_{23}).
\]
Equivalently, after bases have been chosen, the exact ranks of the
product associator defect, the Thom--Sebastiani defect, the
orientation defect, the six projection-idempotence defects, and the
four projection-intertwining defects all vanish.
\end{proposition}

\begin{proof}
The compact Hall product is
\((p_2)_!p_1^!\) on the retained compact-support Hall correspondence.
For three charges, associativity is the equality of the two
correspondences obtained from the stack of two-step extensions. After
choosing finite Borel--Moore bases this equality is exactly the first
displayed matrix identity: the left parenthesisation applies
\(m_{12}\) and then \(m_{12,3}\), while the right parenthesisation
applies \(m_{23}\) and then \(m_{1,23}\).

The vanishing-cycle factors in the two-step correspondence are
identified by Thom--Sebastiani, and the determinant square roots used
for the Joyce--Kontsevich--Soibelman orientation must multiply in the
same way for the two parenthesisations. Those are the two finite
transport equalities \(T_L=T_R\) and \(O_L=O_R\). Compact support is a
separate condition: the product is a compact-support Hall product only
when the retained support summands are actual summands and the four
product maps preserve them. This is precisely the projection
idempotence and projection-intertwining system displayed above.

Every space is finite-dimensional in the fixed height window, so each
condition is an equality of finite rational matrices, hence holds
exactly when the rank of its defect matrix is zero. The finite matrix
witness \texttt{finite\_compact\_hall\_product\_gate} records these
defects and their exact ranks. It checks a supplied compact Hall product
packet; it does not construct the compact-support correspondences, the
global Borel--Moore functor, or the completed compact critical CoHA.
\end{proof}

\begin{proposition}[Finite Hall--Drinfeld double datum criterion]
\label{prop:finite-drinfeld-double-datum-gate}
Fix a finite height window. Let \(Y_H^+\), \(Y_H^-\), \(Y_H^0\), and
\(D_H\) be supplied finite vector spaces intended to be the positive
half, negative half, Cartan part, and finite Hall--Drinfeld double. Let
\[
 G_H\colon Y_H^+\otimes Y_H^-\to \bC((\hbar))
\]
be the reduced Hopf-pairing matrix, let
\[
 \nu_H\colon Y_H^-\otimes Y_H^0\otimes Y_H^+\longrightarrow D_H
\]
be the triangular normal-form matrix, and let
\[
 M_{-+},C_{\mathrm{Dr}}\colon Y_H^-\otimes Y_H^+\longrightarrow D_H
\]
be respectively the supplied mixed product and the matrix obtained from
the Drinfeld cross relation using \(G_H\), the coproduct, the parity
signs, and the Cartan charge action. Suppose also that the finite
coproduct coassociator defect, associator pentagon defect, and centre
compatibility defect have been evaluated as matrices
\[
 \mathcal C_\Delta,\qquad \mathcal P_\Phi,\qquad \mathcal Z_{\mathrm{cent}}.
\]
Then the supplied finite datum is a finite Hall--Drinfeld double datum
exactly when
\[
 \nu_H \text{ is an isomorphism},\qquad
 G_H \text{ is non-degenerate},
\]
\[
 M_{-+}=C_{\mathrm{Dr}},\qquad
 \mathcal C_\Delta=0,\qquad
 \mathcal P_\Phi=0,\qquad
 \mathcal Z_{\mathrm{cent}}=0.
\]
Equivalently, after bases are chosen, the triangular normal-form matrix
is square of full rank, the reduced pairing matrix is square of full
rank, and the exact ranks of the cross-relation, coassociator,
pentagon, and centre defect matrices vanish.
\end{proposition}

\begin{proof}
A Drinfeld double is not determined by its positive half. It is
determined by paired Borel halves, the Cartan part, the cross relation,
and the quasi-Hopf coherence data. In finite dimension the triangular
decomposition is the assertion that multiplication gives a vector-space
normal form
\[
 Y_H^-\otimes Y_H^0\otimes Y_H^+\simeq D_H.
\]
With bases fixed this is exactly the assertion that \(\nu_H\) is a
square full-rank matrix.

The radical quotient has removed precisely the kernel of the Hopf
pairing when \(G_H\) is non-degenerate on the retained positive and
negative halves. This is the square full-rank condition for the reduced
pairing matrix. The mixed product \(Y_H^-\otimes Y_H^+\to D_H\) is then
not an additional choice: it must equal the Drinfeld cross relation
computed from the coproduct, pairing, parity, and Cartan charge. This
is the finite equality \(M_{-+}=C_{\mathrm{Dr}}\).

The remaining three matrices are the finite forms of the double-data
defects named in Definition~\ref{def:compact-oriented-critical-hall-datum}:
coassociativity of the continuous coproduct, the quasi-Hopf pentagon,
and centre compatibility. Each condition is a finite matrix equality,
hence is equivalent to rank zero of its defect. The executable witness
\texttt{finite\_drinfeld\_double\_datum\_gate} records these ranks,
the triangular normal-form rank, and the reduced-pairing rank. It checks
supplied finite double data; it does not construct \(Y_H^-\), \(Y_H^0\),
the coproduct, the completed Hopf pairing, the associator, the centre,
or the global compact Hall--Drinfeld double.
\end{proof}

\begin{proposition}[Finite Rees gluing and exact compact obstructions]
\label{prop:finite-rees-gluing-exact-compact-obstructions}
\label{thm:compact-oriented-critical-hall-cosheaf-criterion}
Let $X=K3\times E$ and fix $(\sigma,S,\mathfrak U)$. For every finite
bound $(N,r,L,m)$ in the bound category $\mathcal I$ of
Definition~\ref{def:compact-oriented-critical-hall-datum}, the Stokes
element
\[
 \Theta_{\hCS\to\Hall}^{\mathrm{Rees},\mathrm{or};N,r,L,m}
 \in
 \mathfrak M_{\hCS,\Hall}^{N,r,L,m}(\mathfrak U)
\]
constructed in Proposition~\ref{prop:gluing-finite-dwr-ran-rees} is a
Maurer--Cartan element and hence a multiplicative natural
transformation from finite hCS observables to the finite oriented Rees
Hall layer on the DWR/Ran nerve.

The passage from these finite Rees maps to a continuous compact
hCS--Hall comparison
\[
 \Theta^{or}_{\hCS\to\Hall}\colon
 \Obs^q_{\hCS}(-,\mathfrak g)
 \longrightarrow
 \CoHA^{or}_{crit}(-)
\]
can be formed, within this compact-passage datum, exactly when
\[
 o_{\mathrm{ML}}=o_{\mathrm{real}}=o_{\mathrm{cpt}}=0.
\]
After this comparison exists, the compact Hall--Drinfeld double
\[
 \mathcal D_\hbar(\CoHA^{or}_{crit}(X))
 =
 (Y^-\widehat\bowtie Y^0\widehat\bowtie Y^+)/
 \operatorname{Rad}\langle-,-\rangle
\]
exists exactly when the double data of
Definition~\ref{def:compact-oriented-critical-hall-datum} are supplied
and
\[
 o_\Delta=o_{\mathrm{pair}}=o_{\mathrm{cent}}=0.
\]
The proposition is a criterion. It proves finite Rees descent; it does
not prove that the compact-passage defects vanish.
\end{proposition}

\begin{proof}
For fixed $(N,r,L,m)$ the finite cyclic critical atlas provides, on
each simplex of the finite DWR/Ran nerve, a relative cyclic critical
model over $\Omega^\bullet(\Delta^p)$ and a face-compatible Rees Hall
map. The relative BV master equation becomes the Maurer--Cartan
equation in the bar--Rees convolution dg Lie algebra. Integration over
$\Delta^p$ converts the de Rham differential on the simplex to the
alternating face differential by Stokes' formula; summing over the
finite nerve gives the displayed finite
$\Theta_{\hCS\to\Hall}^{\mathrm{Rees},\mathrm{or};N,r,L,m}$. The
Thom--Sebastiani rule for disjoint configurations gives
multiplicativity. This is exactly
Proposition~\ref{prop:gluing-finite-dwr-ran-rees}.

The inverse limit over $\mathcal I$ is an exactness question for a
pro-system of complexes. A continuous compact comparison exists only
when the finite maps are compatible under truncation and the
$R^1\!\varprojlim$ term vanishes; this is the
Mittag--Leffler condition $o_{\mathrm{ML}}=0$. If such a continuous
comparison exists, its finite truncations form a Mittag--Leffler system,
so the same condition is forced.

The finite Rees layer is not yet the realised critical CoHA. The target
of the compact comparison is built from equivariant Borel--Moore
homology of vanishing cycles, orientation lines, shifts, and Tate
twists. A monoidal realisation functor compatible with
Thom--Sebastiani and with lci pullback/proper pushforward is therefore
necessary and sufficient to turn the Rees Hall layer into the realised
Hall layer; its failure is precisely $o_{\mathrm{real}}$.

Hall convolution on compact semistable objects is
$(p_2)_!p_1^!$. The shriek pushforward is defined after refinement and
completion only when support remains compact and the required
properness holds for the extension correspondence. This is independent
of the Rees Maurer--Cartan equation: finite Stokes descent supplies a
chain-level map, not compact-support functoriality. Thus
$o_{\mathrm{cpt}}=0$ is the remaining condition for the compact
hCS--Hall comparison.

The Drinfeld double is a construction on paired completed Hopf data.
The realised positive half alone has no Serre-dual negative half, no
Cartan completion, no continuous non-degenerate Hopf pairing, and no
centre compatibility. Supplying these structures and killing
$(o_\Delta,o_{\mathrm{pair}},o_{\mathrm{cent}})$ is exactly the
definition of the displayed completed double. Conversely, any completed
Hall--Drinfeld double contains those data by restriction to its
positive, negative, Cartan, pairing, and centre pieces.
\end{proof}

\begin{corollary}[No shortcut from local critical geometry]
\label{cor:no-shortcut-compact-hall}
None of the following data implies the compact hCS--Hall comparison or
the compact Hall--Drinfeld double on $K3\times E$:
\begin{enumerate}[label=\emph{(\alph*)}, leftmargin=2.5em]
\item the PTVV $(-1)$-shifted symplectic form on the derived moduli
stack;
\item the Brav--Bussi--Joyce local critical charts;
\item the product CY volume form $\Omega_{K3}\wedge\Omega_E$;
\item finite-truncation Kontsevich--Soibelman wall-crossing;
\item local Thom--Sebastiani for one pair of potentials;
\item finite DWR/Ran Rees gluing at every fixed bound $(N,r,L,m)$.
\end{enumerate}
Together these data produce local critical models and finite Rees
descent. They do not imply
$o_{\mathrm{ML}}=o_{\mathrm{real}}=o_{\mathrm{cpt}}=0$, and they do not
supply $(Y^-,Y^0,\Delta_+,\langle-,-\rangle,Z_{\mathrm{cent}})$.
\end{corollary}
\subsection{The $2.5$-obstruction count and the gluing limit}
\label{sec:limit-gluing}

The two-and-a-half-obstruction formulation of
Theorem~\ref{thm:four-obstructions} is structural: $(\mathbf{O_1})$ and
$(\mathbf{O_3})$ are two statements of toric rigidity (absence of a
$T^3$-action on either the total space or on moduli), $(\mathbf{O_4})$
is the categorical shadow of that rigidity (no finite
tilting-quiver), and $(\mathbf{O_2})$ is a Hodge-theoretic obstruction
independent of toric considerations.

\begin{proposition}[Structural logical content]
\label{prop:structural-content}
The four obstructions of Theorem~\ref{thm:four-obstructions} organise
along two axes:
\begin{itemize}[leftmargin=2em]
\item \emph{Toric-rigidity axis.} By Sumihiro 1974 Thm.~1
 (Proposition~\ref{prop:sumihiro-biconditional}),
 $(\mathbf{O_1}) \Leftrightarrow (\mathbf{O_3})$: the absence of a
 $\bC^3$-chart atlas is equivalent to the absence of a rank-$3$ torus
 with isolated fixed points. By Van den Bergh
 \texttt{arXiv:math/0211064} Prop.~3.2.10 plus Bocklandt
 \texttt{arXiv:math/0603558} Thm.~3.1 plus Hartogs
 ($\Gamma(X,\cO_X) = \bC$), the single toric-rigidity proposition
 implies $(\mathbf{O_4})$: no finite Jacobi-algebra tilting generator
 exists on compact $X$.
\item \emph{BCOV axis.} $(\mathbf{O_2})$ is the Atiyah--dilaton
 class $\alpha_{\mathrm{BCOV}}=(\chi_{\mathrm{top}}/24)\,\tr\,
 \mathrm{At}(T_X)\in H^1(X,\cO_X)$ of Costello--Li
 \texttt{arXiv:1606.00365} Prop.~5.2, a single structure-sheaf
 cohomology class independent of any torus or quiver structure.
\end{itemize}
The toric-rigidity axis carries $1.5$ units (one equivalence plus one
further implication); the BCOV axis carries $1$ unit; total $2.5$
units of independent logical content. The four named features
$(\mathbf{C_1})$--$(\mathbf{C_4})$ factor through these two axes.
\end{proposition}

Thus the structural obstruction is two-axis: toric-rigidity and BCOV,
with one further half-step carried by the implication
\((\mathbf{C_2})\Rightarrow(\mathbf{C_3})\).

\begin{remark}[What an extension would look like]
\label{rem:extension-sketch}
A hypothetical extension of chart-wise gluing to compact CY$_3$ would
need to address both axes:
\begin{enumerate}[leftmargin=2em]
\item The toric-rigidity axis: construct a \emph{virtual torus action}
 on moduli via formal stacky enhancements (e.g., Stratum-(iv)
 Borcherds-lift as an $\mathrm{O}(2, n)$-equivariance on a formal
 neighbourhood of the period domain), replacing the toric $T^3$ with
 an automorphic cocycle.
\item The BCOV axis: construct a \emph{BCOV-anomaly-cancelling counter-term}
 in the factorisation-algebra formalism, either through Costello--Li
 obstruction-theoretic means or via a derived-stack enhancement that
 absorbs $\alpha_{\mathrm{BCOV}}$ into the structure cocycle.
\end{enumerate}
Each axis requires non-trivial technology, both axes simultaneously
even more so. The current state is: categorical CoHA exists
(Davison--Meinhardt, stratum-free); chart-wise CoHA does not (obstructed
by both axes).
\end{remark}
\subsection{Chiral-output Booth--Lazarev obstructions for $\Phi^{FA}_3$
on compact CY$_3$}
\label{sec:chiral-output-bl-obstructions}

The chart-gluing obstructions $(\mathbf{O_1})$--$(\mathbf{O_4})$ live on
the \emph{input} side: they measure whether the CY$_3$ category admits a
chart-wise reconstruction. The \emph{output} side of the two-stage
factorisation $\Phi_3 = \mathrm{Sp}^{\mathrm{ch}}_{\Sigma_2, C} \circ
\Phi^{FA}_3$; the chiral factorisation datum
$\Phi^{FA}_3(\Perf\,X)$ sitting on
$\Ran(C) \times \overline{\mathcal{M}}_{g, n}$; carries its own
obstruction triple, orthogonal in content but parallel in structure.
\providecommand{\crv}{\mathrm{crv}}%
\providecommand{\cpt}{\mathrm{cpt}}%
\providecommand{\dgAlg}{\mathrm{dgAlg}}%
\providecommand{\dgCoAlg}{\mathrm{dgCoAlg}}%
The associative Booth--Lazarev Quillen equivalence
$\Omega : \dgCoAlg^{\cpt}_{\crv, k} \rightleftarrows \dgAlg_k : B$
(arXiv:$2304.08409$) leaves three structural upgrades for the chiral
target: Ran-space Smith recognition, genus-tower curvature
compatibility on $\overline{\mathcal{M}}_{g, n}$, and the
sewing-analytic comparison with the Moriwaki IndHilb completion. The
second collapses identically by Mumford's boundary relation; the
residual pair controls whether the chiral output of $\Phi^{FA}_3$ on
a compact CY$_3$ satisfies the Quillen equivalence.

\providecommand{\kch}{\kappa_{\mathrm{ch}}}
\providecommand{\crv}{\mathrm{crv}}
\providecommand{\cpt}{\mathrm{cpt}}
\providecommand{\dgAlg}{\mathrm{dgAlg}}
\providecommand{\dgCoAlg}{\mathrm{dgCoAlg}}
\providecommand{\Dco}{D^{\mathrm{co}}}
\providecommand{\FactCoAlg}{\mathrm{FactCoAlg}}
\providecommand{\FactLie}{\mathrm{FactLie}}
\providecommand{\IndHilb}{\mathrm{IndHilb}}
\providecommand{\Mbar}{\overline{\mathcal{M}}}

\begin{conjecture}[Three chiral-output Booth--Lazarev obstructions]
\label{thm:chiral-output-bl-obstructions}
\ClaimStatusConjectured
Let $X$ be a smooth projective CY$_3$ and
$\Phi^{FA}_3(\Perf\,X)$ the Stage-1 output. For a reference curve
$C \hookrightarrow X$, view the chiral datum as a family over
$\Mbar_{g,n}$ of factorisation coalgebras on $\Ran(C)$; $\Mbar_{g,n}$
is a sewing parameter stack, not a second chiral curve. The existence
of a chiral Booth--Lazarev
Quillen equivalence
\[
\Omega^{\mathrm{ch}} : \FactCoAlg^{\cpt}_{\crv}(\Ran(C))_{\Mbar}
\;\rightleftarrows\;
\FactLie(\Ran(C))_{\Mbar} :B^{\mathrm{ch}}
\]
is obstructed by three cohomological classes:
\begin{enumerate}[label=\emph{$(\mathbf{O}_\arabic*^{\mathrm{BL}})$},
leftmargin=2.5em]
\item \emph{Ran-space Smith recognition:} a class
\[
(\mathbf{O}_1^{\mathrm{BL}}) \;\in\;
H^1\bigl(\Ran(C),\, \mathrm{Aut}_\otimes(\mathrm{Gen}_{\mathrm{BL}})
\bigr)
\]
obstructing factorisation-compatible cofibrant generation of
the coalgebra model structure;
\item \emph{Genus-tower curvature compatibility:} a class
\[
(\mathbf{O}_2^{\mathrm{BL}}) \;\in\;
\bigoplus_{g, n}H^2\bigl(\partial\Mbar_{g, n},\bQ\bigr)
\]
obstructing the clutching-boundary descent
$\partial^* m_0^{(g, n)} = \pi_1^* m_0^{(g_1, n_1+1)} +
\pi_2^* m_0^{(g_2, n_2+1)}$ of the genus-$g$ curvature
$m_0^{(g, n)} \in H^2(\Mbar_{g,n},\bQ)$;
\item \emph{Sewing-analytic IndHilb comparison:} a class
\[
(\mathbf{O}_3^{\mathrm{BL}}) \;\in\;
R^1\!\varprojlim_{\Gamma\in\mathsf G^{\mathrm{st}}_{g,n}}\,
\mathrm{Ext}^1_{\mathrm{top}}\bigl(
{\Dco}^{\mathrm{ch},\Gamma}(\Phi^{FA}_3(\Perf\,X)),\,
\IndHilb^\Gamma(A^{\mathrm{sew}})\bigr)
\]
obstructing the comparison of the algebraic coderived category
with the Moriwaki IndHilb analytic completion over the stable-graph
sewing category.
\end{enumerate}
\end{conjecture}

The obstructions are orthogonal to the chart-gluing obstructions
$(\mathbf{O_1})$--$(\mathbf{O_4})$: they live on the chiral-algebra
output, not on the CY$_3$ input.

\begin{proposition}[$(\mathbf{O}_2^{\mathrm{BL}})$ collapses via Mumford]
\label{prop:O2-BL-mumford-collapse}
The Costello--Li BCOV one-loop curving
(Costello--Li \texttt{arXiv:1606.00365} Prop.~5.2), pulled back along
the $\Ran(C) \times \Mbar_{g, n}$ factorisation tower, takes the form
\[
m_0^{(g, n)} \;=\; \kch \cdot \lambda^{\det}_1(g,n) \;\in\;
H^2(\Mbar_{g,n},\bQ)
\]
with $\kch \in \bQ$ the chain-level Hodge-supertrace invariant and
$\lambda^{\det}_1(g,n)=c_1(\det\mathbb E_{g,n})=c_1(\mathbb E_{g,n})$,
where $\mathbb E_{g,n}=\pi_\ast\omega_\pi$ is the rank-$g$ Hodge
bundle.  The notation $\lambda_i$ is reserved for
$c_i(\mathbb E_{g,n})$; in particular $\lambda_g=c_g(\mathbb E_{g,n})$
is the top Hodge class, not the determinant line. Mumford's boundary
relation
(\emph{Arith.\ \& Geom.}\ II, 1983, Ch.\ III \S 4; Teixidor
\emph{Duke} 56 (1988) Thm.~1.1)
\[
\partial^* \lambda^{\det}_1(g,n) \;=\; \pi_1^* \lambda^{\det}_1(g_1,n_1+1) + \pi_2^*
\lambda^{\det}_1(g_2,n_2+1) \quad (\text{separating, } g = g_1 + g_2),
\qquad
(\partial')^* \lambda^{\det}_1(g,n) \;=\; \lambda^{\det}_1(g-1,n+2)
\quad (\text{non-separating})
\]
forces the clutching-boundary descent identities
$\partial^* m_0^{(g, n)} = \pi_1^* m_0^{(g_1, n_1+1)} + \pi_2^*
m_0^{(g_2, n_2+1)}$ and $(\partial')^* m_0^{(g, n)} = m_0^{(g-1, n+2)}$
to hold identically. Consequently $(\mathbf{O}_2^{\mathrm{BL}}) = 0$
in $H^2(\partial\Mbar_{g, n},\bQ)$: the $\kch$-curving
\emph{is} the Hodge-boundary-descent datum, not a constraint added on
top of it.
\end{proposition}

\begin{proof}
The curving $m_0^{(g, n)} = \kch \cdot \lambda^{\det}_1(g,n)$ is the pullback
of a global rational constant against the Mumford class. Pullback
along the clutching $\partial : \Mbar_{g_1, n_1+1} \times \Mbar_{g_2, n_2+1}
\to \Mbar_{g, n}$ commutes with the scalar $\kch$:
\[
\partial^* m_0^{(g, n)} \;=\; \kch \cdot \partial^* \lambda^{\det}_1(g,n)
\;=\; \kch \cdot (\pi_1^* \lambda^{\det}_1(g_1,n_1+1) + \pi_2^* \lambda^{\det}_1(g_2,n_2+1))
\;=\; \pi_1^* m_0^{(g_1, n_1+1)} + \pi_2^* m_0^{(g_2, n_2+1)}.
\]
Non-separating: $(\partial')^* m_0^{(g, n)} = \kch \cdot (\partial')^*
\lambda^{\det}_1(g,n) = \kch \cdot \lambda^{\det}_1(g-1,n+2) = m_0^{(g-1, n+2)}$. The
identities hold in $\mathrm{Pic}(\partial\Mbar_{g, n})_\bQ$ and in
integral cohomology modulo $2$-torsion. Computer-algebraic verification
at $g \leq 5$ is in Grushevsky--Zakharov \texttt{arXiv:1308.1033} \S 6;
the low-genus sections at $g = 1, 2, 3$ are explicit:
$\partial|_{\delta_0}^* \lambda^{\det}_1(2,0) = \lambda^{\det}_1(1,2)$,
$\partial|_{\delta_1}^* \lambda^{\det}_1(2,0) = 2 \lambda^{\det}_1(1,1)$ (the two
$\pi_i^* \lambda^{\det}_1$ terms, equal by the $\bZ/2$-symmetry of the
$(1, 1)$ split), and similarly at $g = 3$.
\end{proof}

The structural identification deserves emphasis: the one-loop
Costello--Li anomaly and the Mumford boundary relation are not
separately matched; they encode the same Grothendieck--Riemann--Roch
datum viewed through two different descents. The Chern--Gauss--Bonnet
prefactor $\chi(X)/24$ and the Mumford-class weight $1$ both originate
in the Hodge-bundle Quillen determinant $\det R\pi_\ast \omega_{C/S}$,
and the curvature on $\Mbar$ is dictated by the same characteristic
class as the genus-tower curvature on the chiral side.

\begin{remark}[$(\mathbf{O}_1^{\mathrm{BL}})$ is Ran-space, residual]
\label{rem:O1-BL-residual}
The Ran-space cocycle lives in the first \v{C}ech cohomology of the
colimit $\Ran(C) = \mathrm{colim}_n C^n/S_n$ with sheaf coefficients
in $\mathrm{Aut}_\otimes(\mathrm{Gen}_{\mathrm{BL}})$, the sheaf of
tensor-automorphisms of the Booth--Lazarev generating family. The
class measures the failure of factorisation-restriction
$j^*_{n, m}(I_{n+m}) = I_n \boxtimes I_m$ to hold strictly. A
positive-dimensional symmetry can kill the class only after an actual
equivariant lift of the Booth--Lazarev coefficient system has been
constructed and the fixed-locus class has been computed. This is the
$K3\times E$ mechanism on the Heisenberg branch. On strict CY$_3$
(quintic, generic complete intersections) no torus-equivariant
vanishing proof is available because $\Aut^0_s=1$ finite; non-vanishing
still requires a pullback criterion.
\end{remark}

\begin{remark}[$(\mathbf{O}_3^{\mathrm{BL}})$ is sewing-analytic,
class-stratified]
\label{rem:O3-BL-borel}
The chiral OPE on $\Phi^{FA}_3(\Perf\,X)|_{\Ran(C)}$ has pole orders
$\|A_{(k)} B\| \lesssim k!$ in the insertion variable
(Catalan/Gevrey-$1$; cf.\ working\_notes' \v{C}ech-HTT coefficient
convergence, chapters/theory/cy\_to\_chiral.tex Proposition on
class~$\mathbf{M}$ Borel summability). Across a nodal degeneration
$C_t \to C_0$, the asymptotic expansion
$\mu^{\mathrm{ch}}_t = \sum_k (\partial_t^k \mu|_0) \cdot t^k / k!$
has $N$-th truncation error bounded by $(N+1)!^{\kch}$-growth.
Borel summability on the positive ray requires discriminant
$\Delta(Q_\cL) = -256\,\kch^3 \cdot S_4 < 0$, equivalently
$\kch^3 \cdot S_4 > 0$. On class $\mathbf{G}$ (toric CY$_3$ with
terminal shadow $S_2$) the series terminates; on class $\mathbf{L}$
(local $\bP^2$) it grows polynomially; on class $\mathbf{C}$ (K3
Heisenberg--Mukai specialisation with $\kch = 3$, $S_4 > 0$) the
sign-definite regime holds; on class $\mathbf{M}$ (full Gevrey-$1$,
quintic and generic compact CY$_3$) the sign depends on explicit
shadow-tower data. In the non-summable regime
$(\mathbf{O}_3^{\mathrm{BL}})$ persists as an analytic-continuation
class not absorbed by any algebraic construction.
\end{remark}

\begin{proposition}[CFG-type non-compact regime vs.\ compact CY$_3$
residual]
\label{prop:cfg-vs-compact-bl}
The three chiral-output Booth--Lazarev
obstruction classes behave across geometric types as follows:
\[
\begin{array}{|l|c|c|c|}
\hline
\text{Type} & (\mathbf{O}_1^{\mathrm{BL}}) & (\mathbf{O}_2^{\mathrm{BL}}) & (\mathbf{O}_3^{\mathrm{BL}}) \\
\hline
\bR^3 \text{ (topological CS; Costello-Francis-Gwilliam}\ 2026) & 0 & 0 & 0 \\
\bC^3 \text{ (HKR, class } \mathbf{G}) & 0 & 0 & 0 \\
\text{conifold, local }\bP^2 \text{ (classes } \mathbf{G},\,\mathbf{L}) & 0\,(\text{equiv.}) & 0 & 0 \\
	K3 \times E \text{ (Heisenberg, class } \mathbf{C},\, \kch = 3) & 0\,(\text{loc.}) & 0 & 0\,(\text{sign-def.}) \\
	\text{quintic (strict, class } \mathbf{M}) & \text{no equiv. proof} & 0 & \text{Borel-conditional} \\
	\text{generic compact CY$_3$} & \text{criterion needed} & 0 & \text{Borel-conditional} \\
\hline
\end{array}
\]
$(\mathbf{O}_2^{\mathrm{BL}})=0$ is universal by
Proposition~\ref{prop:O2-BL-mumford-collapse}. The
Costello--Francis--Gwilliam $\bR^3$ topological-CS regime and the
$\bC^3$ Dolbeault regime are obstruction-free: the Ran space is
contractible, there is no genus tower, and no compact curve carries
an analytic-continuation anomaly. Compactness removes these automatic
vanishing mechanisms. On a strict compact CY$_3$ without a torus action,
$(\mathbf{O}_1^{\mathrm{BL}})$ must be evaluated by restriction to
explicit curves; the full Gevrey-$1$ shadow tower on class $\mathbf{M}$
leaves $(\mathbf{O}_3^{\mathrm{BL}})$ undecided until the Borel sign is
computed.
\end{proposition}

\begin{proof}[Sketch]
On $\bR^3$: $\Ran(\bR^3)$ is contractible so $H^1(\Ran(\bR^3),
\cdot) = 0$; no $\Mbar$, no degeneration. On $\bC^3$: $\Ran(\bC)$
contractible and non-compact (no $H^1$ in relevant coefficients),
$\alpha_{\mathrm{BCOV}}(\bC^3) = 0$ (Costello--Li target $H^1(\bC^3,
\cO) = 0$), and the class-$\mathbf{G}$ shadow tower terminates.
On $K3 \times E$ with the Heisenberg specialisation, the
$E$-translation torus $T = E$ acts on $\Ran(E)$ by translation,
$\Ran(E)^T = $ empty generically so equivariant cohomology localises
to the fixed fibre (after rationalisation); the Dunn--Lurie descent
comparison on \(\Ran(C)\times\Mbar\) lifts
the equivariant class to a $\bQ$-cycle that vanishes upon
rationalisation. On strict compact CY$_3$, $\Aut^0_s$ finite forecloses
this particular localisation proof; it does not by itself compute the
Smith class.
\end{proof}

\begin{remark}[Four-$\kappa_\bullet$ discipline in the curving]
The curving is $m_0^{(g, n)} = \kch \cdot \lambda^{\det}_1(g,n)$, not
$\kcat \cdot \lambda^{\det}_1(g,n)$, nor
$\kBKM \cdot \lambda^{\det}_1(g,n)$, nor
$\kfib \cdot \lambda^{\det}_1(g,n)$. The four invariants differ on $K3 \times E$:
$\kch = 0$ (Dolbeault Hodge-supertrace on total space) versus
$\kch^{\mathrm{Heis}} = 3$ (after Stage-2 chiral Heisenberg--Mukai
specialisation) versus $\kcat = 0$ (Künneth-multiplicative) versus
$\kBKM(\Delta_5) = 5$ (Borcherds) versus $\kfib = 24$
(Mukai-lattice fibre). The curving of the chiral Booth--Lazarev
genus-tower uses the Hodge-supertrace invariant $\kch$, sharpened
by the Stage-2 specialisation if the output is the Heisenberg
branch.
\end{remark}

\begin{remark}[Stage-2 specialisation absorbs part of
$(\mathbf{O}_1^{\mathrm{BL}})$]
The specialisation
$\mathrm{Sp}^{\mathrm{ch}}_{\Sigma_2, C} : \mathrm{FactAlg}_3
\to \mathrm{FactAlg}_{(E_1, C)}$ restricts the factorisation
algebra to a specific curve $C$. The \v{C}ech-Ran cocycle
representing $(\mathbf{O}_1^{\mathrm{BL}})$ on $\Ran(X) \supset \Ran(C)$
pulls back along $j^*_{C/X}$; the restriction quotients away the
higher-$n$ stratum cocycles that do not restrict to $C^n$. The
residual class lies in the image of
$H^1(\Ran(X), \mathrm{Aut}_\otimes) \to H^1(\Ran(C),
\mathrm{Aut}_\otimes)$, strictly smaller than the starting class.
The sewing-analytic obstruction $(\mathbf{O}_3^{\mathrm{BL}})$ is
intrinsic to $C$'s $\Mbar_{g, n}$-parameter space and unaffected
by the Stage-2 transverse choice.
\end{remark}

\begin{remark}[Residual obstruction count]
\label{rem:bl-obstruction-count}
Among the three chiral-output Booth--Lazarev obstructions, the
genus-tower obstruction $(\mathbf{O}_2^{\mathrm{BL}})$ collapses
identically by Mumford's boundary relation: it is not an
independent obstruction, but rather the structural identification
between the one-loop BCOV curving $\chi(X)/24$ and the Hodge
determinant bundle $\lambda_g$. The residual pair
$(\mathbf{O}_1^{\mathrm{BL}})$, $(\mathbf{O}_3^{\mathrm{BL}})$
carries the full obstruction content. The count matches the two-axis
structural content of Proposition~\ref{prop:structural-content} on the
chart-gluing side ($\mathbf{O}_2^{\mathrm{BL}}$ collapse parallels
the categorical implication $(\mathbf{O_1}) \Rightarrow
(\mathbf{O_4})$ on the chart side, leaving two irreducible axes per
side). The chart-side axes are toric-rigidity and BCOV; the
output-side axes are Ran-space Smith recognition and sewing-analytic
comparison. The two pairs are logically independent: compact CY$_3$
can be chart-obstructed and output-obstructed simultaneously, or
chart-obstructed but output-equivariant-trivial (e.g., $K3 \times E$
on the Heisenberg branch).
\end{remark}
\subsubsection*{Explicit chain-level representatives for the residual
pair}
\label{sec:residual-pair-chain-level}

The residual obstructions $(\mathbf{O}_1^{\mathrm{BL}})$ and
$(\mathbf{O}_3^{\mathrm{BL}})$ admit concrete chain-level
representatives as parallel Mittag--Leffler defects in two orthogonal
directions of the bi-filtered space $\Ran(C) \times \Mbar$. The
$(\mathbf{O}_1^{\mathrm{BL}})$-class is a
stratum-count Mittag--Leffler defect; the
$(\mathbf{O}_3^{\mathrm{BL}})$-class is a genus-tower Mittag--Leffler
defect. The Mumford collapse of $(\mathbf{O}_2^{\mathrm{BL}})$ is the
\emph{curvature}-compatibility identification between the two
directions; the two identity-type residuals remain.

\begin{proposition}[$(\mathbf{O}_1^{\mathrm{BL}})$ as \v{C}ech cocycle]
\label{prop:O1-BL-cech-cocycle}
Fix a smooth reference curve
$C \hookrightarrow X$ and let $V_n = \{(x_1, \dots, x_n) \in C^n/S_n :
x_i \neq x_j\}$ stratify the Ran space. The cocycle
\[
\sigma_{\mathrm{Smith}}^{(n, m)} \;:=\;
[j^*_{n, m}\,I_{n+m}\, -\, I_n \boxtimes I_m] \;\in\;
\Aut_\otimes(\Gen_{\mathrm{BL}})(V_n \cap V_m)
\]
on the open locus of distinct $(n+m)$-tuples satisfies the
factorisation-associativity triple-overlap cocycle condition
\[
\sigma^{(n, m+k)} - \sigma^{(n+m, k)} + \sigma^{(n, m)}
 - \sigma^{(m, k)} = 0
\qquad \text{in } V_n \cap V_m \cap V_k,
\]
and its cohomology class realises the obstruction
\[
(\mathbf{O}_1^{\mathrm{BL}}) = [\sigma_{\mathrm{Smith}}]
\in H^1(\Ran(C), \Aut_\otimes(\Gen_{\mathrm{BL}}))
\]
in the finite-set/Ran descent diagram.  It becomes a sequential
$\varprojlim^1$ only after a cofinal diagonal subsystem and compatible
transition maps have been chosen; no genus- or point-count tower is
part of the definition.
Pulled back along an acyclic analytic chart cover
$\{U_\alpha\}$ of $X$, the cocycle is represented by a
Dolbeault--\v{C}ech class
$\tilde\sigma_{\alpha\beta\gamma} \in
\Gamma(U_{\alpha\beta\gamma}, \Omega^{0,1} \otimes
\Aut_\otimes^{\mathrm{FM}})$ on triple overlaps, with
$\Aut_\otimes^{\mathrm{FM}}$ the Bondal--Orlov Fourier--Mukai
autoequivalence sheaf.
\end{proposition}

\begin{proof}
The factorisation tensor structure on
$\FactCoAlg^{\cpt}_{\crv}(\Ran(C))$ specifies that the generating
cofibrations on distinct configurations satisfy
$j^*_{n, m}(I_{n+m}) = I_n \boxtimes I_m$ strictly or up to
automorphism of the generating family. The discrepancy is the
$1$-cochain above. Associativity of the factorisation tensor
(Beilinson--Drinfeld Ch.~$3$ Def.~$3.4.9$) gives the cocycle
condition. By Ayala--Francis \texttt{arXiv:1206.5522} Thm.~$1.1$,
$H^\bullet(\Ran(C), \cF)$ is computed by the
Bousfield--Kan spectral sequence on the \v{C}ech nerve of the finite-set
diagram; a sequential Mittag--Leffler description is a further
cofinality theorem, not part of the Ran geometry itself. The chart-cover Dolbeault pull-back
uses Cartan's Theorem~B acyclicity on $U_{\alpha\beta\gamma}
\simeq \bC^3$ intersections and the Bondal--Orlov uniqueness of the
Fourier--Mukai kernel up to shift/line-bundle twist
(Bondal--Orlov \texttt{arXiv:alg-geom/9712029}). The Kapranov
$L_\infty$-structure \texttt{arXiv:math/9811023} Thm.~$2.7$
identifies the Atiyah-class component of the triple-overlap
discrepancy with the $\Omega^{0, 1}$-Dolbeault cocycle.
\end{proof}

\begin{proposition}[$(\mathbf{O}_3^{\mathrm{BL}})$ as stable-graph sewing defect]
\label{prop:O3-BL-ml-defect}
This statement is made on the stable-graph sewing diagram with residue
maps along new half-edges. There is no canonical genus-only inverse
tower $\Mbar_{g+1}\to\Mbar_g$.
Let $\mathsf G^{\mathrm{st}}_{g,n}$ be the stable-graph category of
boundary strata of $\Mbar_{g,n}$.  The sewing-completed state spaces of
$\cA = \PhiFA_3(\Perf\,X)|_{\Ran(C)}$, parameterised over $\Mbar$, form
a diagram
$\Gamma\mapsto A^{\mathrm{sew}}_\Gamma(\cA)$ under clutching pullback
and residue along the new half-edges.  Its derived inverse-limit defect is
\[
(\mathbf{O}_3^{\mathrm{BL}}) = [\sigma_{\mathrm{sew}}]
\;\in\; R^1\!\varprojlim_{\Gamma\in\mathsf G^{\mathrm{st}}_{g,n}}\,
A^{\mathrm{sew}}_\Gamma(\cA).
\]
The chain-level representative is the leading non-nuclear
contribution to the chiral OPE asymptotic across a nodal
degeneration $C_t \to C_0$, explicitly
\[
\sigma_{\mathrm{sew}}^{(N)}(A, B; z, t)
\;=\;
\sum_{k=0}^{N} \frac{t^k}{z^{k+1}}\,
\mathrm{Res}_{\mathrm{node}}^{(k)}(A, B) \;-\;
\mu^{\mathrm{ch}}_t(A(z_1), B(z_2)),
\]
with truncation-error bound
$\|\sigma_{\mathrm{sew}}^{(N)}\| \lesssim t^{N+1} \cdot
(N+1)!^{\kch}$ (Gevrey-$1$). The Borel-summability criterion
$\Delta(Q_\cL) = -256\,\kch^3 S_4 < 0$
(equivalently $\kch^3 S_4 > 0$) controls whether $N \to \infty$
converges to zero in the nuclear IndHilb topology after the stable-graph
sewing diagram has been constructed.
\end{proposition}

\begin{proof}
The boundary of $\Mbar_{g,n}$ is indexed by stable graphs, not by a
linear genus tower.  Separating clutching maps have product source
$\Mbar_{g_1,n_1+1}\times\Mbar_{g_2,n_2+1}$; non-separating clutching
has source $\Mbar_{g-1,n+2}/S_2$.  The sewing diagram therefore carries
both pullback along clutching and residue/trace along the two new
markings.
The chiral OPE Taylor-expands in the sewing parameter with
Gevrey-$1$ coefficient growth (Vol~I $\mathbf{MC5}$ sewing problem;
cy\_to\_chiral.tex Proposition on class $\mathbf{M}$ Borel
summability). The Mittag--Leffler defect of a Gevrey-$1$ tower
is controlled by the Borel discriminant of its shadow-tower
quadratic $Q_\cL$; the identity $\Delta(Q_\cL) = -256 \kch^3 S_4$
is Proposition~$3.6$ of cy\_to\_chiral.tex, the shadow-tower
$\alpha^2$-cancellation. Sign-definite negative $\Delta$ gives
complex-conjugate Borel-plane singularities off the positive ray,
Borel summability, and hence $\varprojlim^1_g = 0$ by the
standard surjection-implies-$\varprojlim^1$-vanishing
(Hartshorne \emph{Residues and Duality} Appendix A) only on a
separately constructed non-separating residue subsystem.  The full
Booth--Lazarev obstruction is the stable-graph derived-limit class
above.
\end{proof}

\begin{remark}[Parallel Mittag--Leffler structure]
\label{rem:residual-ml-parallel}
Both residual obstructions are descent defects in orthogonal
filtration directions of $\Ran(C) \times \Mbar$: the Ran-space class
is indexed by finite sets, surjections, and diagonal restrictions,
while the sewing class is indexed by stable graphs. The
Mumford collapse of $(\mathbf{O}_2^{\mathrm{BL}})$ is the
\emph{curvature}-compatibility identification across this bi-filtration
(the $\kch$-curving \emph{is} the Hodge-boundary-descent datum).
The two identity-type defects remain independent: neither controls
the other, and both must vanish for the full chain-level Quillen
equivalence on $\Ran(C) \times \Mbar$.
\end{remark}

\begin{theorem}[Residual-pair values on three canonical CY$_3$]
\label{thm:residual-pair-three-values}
The $K3\times E$ entry uses rational equivariant localisation and the
sign condition $S_4>0$. The quintic entry gives a detection criterion,
not an unconditional non-vanishing theorem.
The residual pair
$(\mathbf{O}_1^{\mathrm{BL}}), (\mathbf{O}_3^{\mathrm{BL}})$
on the Stage-1 output $\PhiFA_3(\Perf\,X)$ admits the following
first-principles values:
\begin{enumerate}[label=\emph{(\roman*)}, leftmargin=2.5em]
\item $X = \bC^3$: both vanish trivially.
$(\mathbf{O}_1^{\mathrm{BL}})(\bC^3) = 0$ because
$\Ran(\bC^3)$ is contractible (Beilinson--Drinfeld Ch.~$3$
Lemma~$3.4.4$) and $H^{>0}(\Ran(\bC^3), \cdot) = 0$.
$(\mathbf{O}_3^{\mathrm{BL}})(\bC^3) = 0$ because $\bC^3$ is
non-compact with no relevant genus tower: the sewing-completion
system $\{A^{\mathrm{sew}, g}(\CoHA(\bC^3))\}$ is constant in $g$
and Mittag--Leffler trivially.

\item $X = K3 \times E$ on the Heisenberg--Mukai specialisation branch:
both residuals vanish after two special-geometry checks. The class
$(\mathbf{O}_1^{\mathrm{BL}})$ vanishes rationally when the
$T=E$-translation localisation map identifies
$[\sigma_{\mathrm{Smith}}]\otimes\bQ$ with the fixed-locus class. The
class $(\mathbf{O}_3^{\mathrm{BL}})$ vanishes when
$\kch^{\mathrm{Heis}}(K3 \times E)=3$ and
$S_4(\cH_{\mathrm{Muk}})>0$, so that
$\Delta(Q_\cL)=-256\cdot 27\cdot S_4<0$.

\item $X = Q \subset \bP^4$ the generic quintic threefold: the
stratum-free data do not supply a vanishing mechanism. Matsumura--Oguiso
finiteness gives $\Aut^0(Q)=1$, so the equivariant localisation proof
available on $K3\times E$ has no analogue. A rigid line
$\ell\subset Q$ gives a detection criterion for
$(\mathbf{O}_1^{\mathrm{BL}})$: the class is non-zero if
\[
 \ell^*[\sigma_{\mathrm{Smith}}]\;\neq\;0
 \quad\text{in}\quad
 H^1(\Ran(\ell),\Aut_\otimes(\Gen_{\mathrm{BL}})).
\]
The sewing class $(\mathbf{O}_3^{\mathrm{BL}})(Q)$ is Borel-conditional
on the sign of $S_4(Q)$: the tree-level Hodge supertrace
$\kch(Q) = \chi(\cO_Q) = 0$ makes the naive discriminant
degenerate, and the BV-corrected invariant
$a_{\mathrm{BV}}(Q) = \chi(Q)/24 = -25/3$ is the correct one-loop
input entering the Costello--Li curving, giving
$\Delta(Q_\cL)(Q) = 32000\, S_4(Q)/27$; sign-definiteness
requires $S_4(Q) < 0$, pending explicit Bershadsky--Cecotti--Ooguri--Vafa
recursion.
\end{enumerate}
\end{theorem}

\begin{proof}
$(\bC^3)$ Contractibility of $\Ran(\bC^3)$ follows from the uniform
scaling contraction $(x_1, \dots, x_n) \mapsto (t \cdot x_1, \dots,
t \cdot x_n)$, $t \in [0, 1]$: for $t > 0$ the scaled configuration
has distinct points (since distinct $x_i \neq x_j$ force distinct
$t x_i \neq t x_j$), so the path lies entirely in
$\Ran(\bC^3) \setminus \Delta$ until it terminates at the apex at
$t = 0$. Ayala--Francis \texttt{arXiv:1206.5522} Thm.~$1.1$ extends
this to the descent of $H^\bullet(\Ran(\bC^3), \cF)$ for any
coefficient sheaf. The sewing-completion on $\CoHA(\bC^3) =
Y^+(\widehat{\mathfrak{gl}}_1)$ (Schiffmann--Vasserot
\texttt{arXiv:1202.2756} Thm.~$1.1$) is genus-independent (no
compact stable curves map to non-compact $\bC^3$).

$(K3 \times E)$ The translation torus $T = E$ acts on $\Ran(E)$ by
$(y_1, \dots, y_n) \mapsto (y_1 + \tau, \dots, y_n + \tau)$ for
$\tau \in E$. For generic (non-torsion) $\tau$, the only fixed
configuration is the empty one; $\Ran(E)^T = \emptyset$ apart from
the base point. If the Smith cocycle is represented in the
$T$-equivariant Booth--Lazarev coefficient sheaf, Atiyah--Bott
localisation identifies its rational class with the fixed-locus class,
which is zero. This proves rational vanishing exactly under the
equivariant-lift hypothesis stated in the theorem. For the sewing
obstruction:
the Heisenberg--Mukai specialisation has $\kch^{\mathrm{Heis}} = 3$
(Stage-$2$ specialisation of $\Phi^{FA}_3$; cy\_to\_chiral.tex
Theorem on $K3 \times E$ Heisenberg) and $S_4 > 0$ from the
Mukai pairing's positive-cone contribution (Gritsenko--Nikulin
$2002$ Adv. Math. 165 Prop.~$2.5$, positivity of Fourier
coefficients on the positive cone). Both inputs combined give the
sign-definite discriminant and Borel summability.

$(Q)$ Matsumura's theorem: for a smooth hypersurface $X_d \subset
\bP^n$ of degree $d \geq n + 2$, $\Aut(X_d)$ is finite. The quintic
has $d = 5 = n + 1$, borderline; Oguiso $2005$ Sci. China Math. 48
extends to the borderline case for generic $Q$ with
$\mathrm{Pic}(Q) = \bZ \cdot H$, giving $\Aut^0(Q) = 1$. No algebraic
torus acts, so the $K3\times E$ localisation proof cannot be repeated.
On a rigid rational curve $\ell \subset Q$ (Harris $1979$;
$N_{\ell/Q} = \cO(-1) \oplus \cO(-1)$), the tangent bundle
$T_\ell=\cO_\ell(2)$ has $c_1(T_\ell)=2$, so
$\At(T_\ell)\in H^1(\ell,\Omega^1_\ell)\simeq\bC$ is non-zero. The
pullback $\ell^*[\sigma_{\mathrm{Smith}}]$ is therefore the correct
chain-level criterion for non-vanishing; absence of torus equivariance alone
is only an obstruction to the known vanishing proof.

For the quintic $(\mathbf{O}_3^{\mathrm{BL}})$: the Hodge diamond of
$Q$ has $h^{0, 0} = 1, h^{0, 1} = 0, h^{0, 2} = 0, h^{0, 3} = 1$
(Griffiths--Harris on quintic threefolds), so
$\kch(Q) = 1 - 0 + 0 - 1 = 0$. The naive Borel discriminant
$\Delta = -256 \cdot 0 \cdot S_4 = 0$ is degenerate, so the
Proposition~$3.6$ of cy\_to\_chiral.tex does not apply directly.
The BV anomaly coefficient
$a_{\mathrm{BV}}(Q) = \zeta(0)\big|_{L^5} \cdot \mathrm{Res}_{\bar\partial}
\alpha_{\mathrm{BCOV}}(Q) = \chi(Q)/24 = -25/3$
(Remark~\ref{rem:kappa-ch-bv-distinction}) enters the one-loop
BCOV curving; substituting $a_{\mathrm{BV}}$ for $\kch$ in the
discriminant formula gives
$\Delta = -256 \cdot (-125/27) \cdot S_4 = 32000\, S_4/27$,
sign-definite negative iff $S_4(Q) < 0$.
\end{proof}

\begin{remark}[Sharpening: $\kch$ vs $a_{\mathrm{BV}}$ on the quintic]
\label{rem:kch-vs-kchBV-quintic-sharpened}
The distinction between tree-level Hodge supertrace $\kch$ and
BV one-loop anomaly coefficient $a_{\mathrm{BV}}$ is load-bearing on the
quintic: $\kch(Q) = 0$ from the Hodge diamond (trivial
$h^{0,1}, h^{0,2}$), but $a_{\mathrm{BV}}(Q) = \chi(Q)/24 = -25/3$ from the
topological Euler-characteristic-weighted one-loop trace. The
Borel-summability criterion for $(\mathbf{O}_3^{\mathrm{BL}})$ on
the quintic uses $a_{\mathrm{BV}}$, not $\kch$, because the sewing-obstruction
tower is sensitive to the BV-regularised chiral-trace residue, not
to the tree-level Dolbeault supertrace.

The four scalar invariants attached to the quintic, with the BV coefficient
kept separate, are
\[
\kch(Q) = 0, \quad \kcat(Q) = 0, \quad \kBKM \text{ is not defined for } Q,
\quad \kfib(Q) \text{ is not defined},
\quad a_{\mathrm{BV}}(Q) = -25/3.
\]
The Borel-summability criterion uses the BV coefficient
precisely because the tree-level Hodge supertrace vanishes and would
otherwise make the criterion degenerate.
\end{remark}

\begin{corollary}[Residual-pair criterion, not a compact-CY$_3$
classification]
\label{cor:residual-pair-stratification}
The residual pair
$(\mathbf{O}_1^{\mathrm{BL}}), (\mathbf{O}_3^{\mathrm{BL}})$
is controlled by two independent vanishing criteria on a smooth
projective CY$_3$:
\begin{itemize}[leftmargin=2em]
\item $(\mathbf{O}_1^{\mathrm{BL}})=0$ follows from an established
factorisation-compatible acyclic Ran cover, or from an equivariant
localisation theorem that identifies the rational Smith cocycle with a
vanishing fixed-locus class. If no algebraic torus acts, this proof
mechanism is absent; non-vanishing still requires a pullback criterion such
as the rigid-line criterion in
Theorem~\ref{thm:residual-pair-three-values}.
\item $(\mathbf{O}_3^{\mathrm{BL}})=0$ follows from the sign-definite
Borel criterion $\Delta(Q_\cL)=-256\,\kch^3S_4<0$ after the correct
chiral trace invariant has been chosen. If the discriminant is
degenerate or its sign is unknown, the sewing obstruction is
Borel-conditional and must be computed case by case.
\end{itemize}
The $K3\times E$ Heisenberg--Mukai branch satisfies both criteria under
the localisation and positivity hypotheses listed in
Theorem~\ref{thm:residual-pair-three-values}. A strict generic compact
CY$_3$ supplies neither criterion from the stratum-free data alone.
\end{corollary}

\begin{proof}
The first criterion is the definition of the Smith cocycle as a
\v{C}ech--Ran class: either the coefficient system is acyclic on the
chosen factorisation cover, or an equivariant localisation theorem
computes the class on the fixed locus. These are sufficient conditions,
not a classification of all possible vanishing mechanisms. The second
criterion is Proposition~\ref{prop:O3-BL-ml-defect}: a negative Borel
discriminant places the singularities away from the positive ray and
kills the Mittag--Leffler defect. If the sign is unknown, the theorem
does not decide the class. The $K3\times E$ and quintic statements are
the two worked applications of these criteria in
Theorem~\ref{thm:residual-pair-three-values}.
\end{proof}
\subsection{Summary}
\label{sec:summary}

The results of this part:
\begin{enumerate}[leftmargin=2em]
\item Four named obstructions $(\mathbf{O_1})$--$(\mathbf{O_4})$
 to chart-wise gluing on a compact CY$_3$
 (Theorem~\ref{thm:four-obstructions}), each derived from a primary
 source: Demazure--Fulton \S3.4 for $(\mathbf{O_1})$; Costello--Li
 \texttt{arXiv:1606.00365} Prop.~5.2 for $(\mathbf{O_2})$; Bogomolov
 1974 + Beauville 1983 JDG 18 \S2 + Matsumura 1963 Proc Japan Acad 39
 for $(\mathbf{O_3})$; Van den Bergh \texttt{arXiv:math/0211064} +
 Bocklandt \texttt{arXiv:math/0603558} Thm.~3.1 + Hartogs for
 $(\mathbf{O_4})$.
\item The Sumihiro 1974 biconditional $(\mathbf{O_1}) \Leftrightarrow
 (\mathbf{O_3})$ (Proposition~\ref{prop:sumihiro-biconditional}) and
 the derived implication $(\mathbf{O_1}) \Rightarrow (\mathbf{O_4})$,
 collapsing the four-feature enumeration to $2.5$ units of independent
 logical content (Proposition~\ref{prop:structural-content}).
\item Compact CY$_3$ examples (quintic, bicubic, Schoen) that lack
 $K3 \times E$'s exceptional triple alignment and hence admit no
 chart-wise assembly (Section~\ref{sec:compact-no-fibration}).
\item The stratum-free receptacles (Davison--Meinhardt critical-Hall
 target, Kontsevich--Soibelman motivic wall-crossing, two-stage
 factorisation target) that can be formulated for compact CY$_3$ only
 with their own orientation, realisation, convergence, and framing
 hypotheses
 (Proposition~\ref{prop:stratum-free}) and the chart-wise machinery
 (Schiffmann--Vasserot stalks, GLY bond factors, shuffle algebras,
 finite quivers) that does not
 (Proposition~\ref{prop:transfer-failures}).
\item The trace-splitting obstruction to a global NCCR on
 \(K3\times E\), the four failed construction routes, and the
 Serre-equivariant quasi-NCCR substitute
 (Proposition~\ref{prop:k3e-five-obstructions};
 Remark~\ref{rem:serre-equivariant}).
\item The finite DWR/Ran Rees hCS--Hall gluing
 (Proposition~\ref{prop:finite-rees-gluing-exact-compact-obstructions}):
 for every finite bound $(N,r,L,m)$ the Stokes element is a
 Maurer--Cartan element on the finite nerve. The compact comparison
 requires exactly
 $o_{\mathrm{ML}}=o_{\mathrm{real}}=o_{\mathrm{cpt}}=0$, and the
 compact Hall--Drinfeld double additionally requires
 $o_\Delta=o_{\mathrm{pair}}=o_{\mathrm{cent}}=0$
 (Corollary~\ref{cor:no-shortcut-compact-hall}).
\item Recovery statement (Theorem~\ref{thm:recovery-summary}) and
 no-stratum-free-$G(X)$ theorem (Theorem~\ref{thm:no-stratum-free-gx}):
 the critical-Hall/BPS package is conditional, MMNS motivic generating
 functions and numerical BPS invariants survive, and the
 Borcherds--Gritsenko automorphic formula is available only in
 stratum~(iv). Chart-wise stalks, bond factors, shuffles, quivers,
 compact framed $\Phi_3$, global $G(X)$, Hall--Drinfeld doubles, and
 the automorphic cocycle outside stratum~(iv) are not supplied by
 stratum-free compact data.
\item Three chiral-output Booth--Lazarev obstructions
 $(\mathbf{O}_1^{\mathrm{BL}}), (\mathbf{O}_2^{\mathrm{BL}}),
 (\mathbf{O}_3^{\mathrm{BL}})$ on the Stage-1 output side
 (Theorem~\ref{thm:chiral-output-bl-obstructions}), with
 $(\mathbf{O}_2^{\mathrm{BL}})$ collapsing identically via Mumford's
 boundary relation (Proposition~\ref{prop:O2-BL-mumford-collapse})
 and residual pair $(\mathbf{O}_1^{\mathrm{BL}})$,
 $(\mathbf{O}_3^{\mathrm{BL}})$ controlling whether chiral
 bar-cobar closes on $\Ran(C) \times \Mbar$.
\item Explicit chain-level representatives for the residual pair:
 $(\mathbf{O}_1^{\mathrm{BL}})$ as the
 factorisation-associativity \v{C}ech cocycle
 $[\sigma_{\mathrm{Smith}}] = [j^*_{n,m} I_{n+m} - I_n \boxtimes I_m]$,
 with a sequential Mittag--Leffler description only after a cofinal
 finite-set subsystem and compatible transition maps have been chosen
 (Proposition~\ref{prop:O1-BL-cech-cocycle}); and
 $(\mathbf{O}_3^{\mathrm{BL}})$ as the stable-graph sewing
 derived-limit defect represented by the leading non-nuclear term in
 the chiral OPE asymptotic across a nodal degeneration
 (Proposition~\ref{prop:O3-BL-ml-defect}), with
 Borel-summability criterion $\Delta(Q_\cL) = -256\,\kch^3 S_4 < 0$.
\item Three-example residual-pair values
 (Theorem~\ref{thm:residual-pair-three-values}): both vanish on
 $\bC^3$ (contractible Ran, no genus tower); both vanish on the
 $K3 \times E$ Heisenberg branch under the explicit localisation and
 sign-definite Borel hypotheses; the quintic lacks the equivariant
 vanishing proof for $O_1$ and leaves $O_3$ Borel-conditional on the
 sign of $S_4(Q)$, using the BV-corrected $a_{\mathrm{BV}}(Q) = -25/3$ because
 tree-level $\kch(Q)=0$ is degenerate.
\item Residual-pair criterion
 (Corollary~\ref{cor:residual-pair-stratification}): $O_1$ vanishes by
 acyclic Ran descent or verified equivariant localisation; $O_3$
 vanishes by a sign-definite Borel discriminant. The corollary is a
 criterion, not a classification of all compact CY$_3$'s.
\end{enumerate}

The chart-side gluing machinery of Parts~II--VIII is sharp: it covers
$K3 \times E$ and the toric CY$_3$ family without compact $4$-cycles,
and no further. The boundary is dictated by the two axes
(toric-rigidity + BCOV); any extension beyond requires engaging both
axes simultaneously with non-classical machinery. The output-side
Booth--Lazarev obstructions form a parallel two-axis structure after
the Mumford collapse of $(\mathbf{O}_2^{\mathrm{BL}})$: Ran-space
Smith recognition $(\mathbf{O}_1^{\mathrm{BL}})$ and sewing-analytic
IndHilb comparison $(\mathbf{O}_3^{\mathrm{BL}})$, logically
independent of the chart-side axes. On the Heisenberg-specialised
branch of $K3 \times E$ both output obstructions vanish (torus
equivariance + sign-definite Borel summability); on strict compact
CY$_3$ the stratum-free data alone do not supply either vanishing
criterion.
